\documentclass{article}

% if you need to pass options to natbib, use, e.g.:
%     \PassOptionsToPackage{numbers, compress}{natbib}
% before loading neurips_data_2022

% ready for submission
% \usepackage{neurips_data_2023}

% to compile a preprint version, add the [preprint] option, e.g.:
% \usepackage[preprint]{neurips_data_2023}
% This will indicate that the work is currently under review.

% to compile a camera-ready version, add the [final] option, e.g.:
%     \usepackage[final]{neurips_data_2022}

% to avoid loading the natbib package, add option nonatbib:
   \usepackage[nonatbib]{neurips_data_2023}

% Submissions to the datasets and benchmarks are typically non anonymous,
% but anonymous submissions are allowed. If you feel that you must submit 
% anonymously, you can compile an anonymous version by adding the [anonymous] 
% option, e.g.:
%     \usepackage[anonymous]{neurips_data_2022}
% This will hide all author names.

% \usepackage{wrapfig}
\usepackage[utf8]{inputenc} % allow utf-8 input
\usepackage[T1]{fontenc}    % use 8-bit T1 fonts
% \usepackage{hyperref}       % hyperlinks
\usepackage{url}            % simple URL typesetting
\usepackage{booktabs}       % professional-quality tables
\usepackage{amsfonts}       % blackboard math symbols
\usepackage{nicefrac}       % compact symbols for 1/2, etc.
\usepackage{microtype}      % microtypography
\usepackage{xcolor}         % colors
\usepackage{listings}
\usepackage{epsfig}
\usepackage{graphicx}
\usepackage{amsmath}
\usepackage{amssymb}
\usepackage{multirow}
\usepackage{color, colortbl}
\usepackage{subcaption}
\usepackage{pifont}
\usepackage{todonotes}
\usepackage{bbding}
\usepackage{wasysym}
\usepackage{comment}
\usepackage{bm}
\usepackage{makecell}
\usepackage[misc]{ifsym}

\definecolor{TRAD_EVAL}{HTML}{edede9}
\definecolor{NVL_EVAL}{HTML}{edede9}

% \definecolor{TRAD_EVAL}{HTML}{E6ECE3}
% \definecolor{NVL_EVAL}{HTML}{ffefe0}

\newcommand\blfootnote[1]{%
  \begingroup
  \renewcommand\thefootnote{}\footnote{#1}%
  \addtocounter{footnote}{-1}%
  \endgroup
}


\lstset{                         
showstringspaces=false,                
frame=single,                           
}

\definecolor{citecolor}{HTML}{0071BC}
\definecolor{linkcolor}{HTML}{ED1C24}
\definecolor{NEED_TRAIN}{HTML}{ffefe0}
\newcommand{\subfig}[1]{\textcolor{linkcolor}{-#1}}

\newcommand{\jingkang}[1]{{\color{pink}[Jingkang: #1]}}
\newcommand{\jack}[1]{{\color{red}[Jack: #1]}}

\makeatletter
\renewcommand\paragraph{
  \@startsection{paragraph} % name
  {4} % level
  {\z@} % indent
  {.5em \@plus1ex \@minus.2ex} % beforeskip
  {-1.5em} % afterskip
  {\normalfont\normalsize\bfseries} % style
}

\usepackage[pagebackref=true, breaklinks=true, colorlinks,citecolor=citecolor, linkcolor=linkcolor, bookmarks=false]{hyperref}

% \title{FunQA: A Benchmark for Unraveling Video Comprehension with Surprising Clips}
\begin{document}
% \title{
% \begin{wrapfigure}{}{}
%     \raggedright
%     \includegraphics[width=0.044\textwidth]{figures/logo.png}
% \end{wrapfigure}
% FunQA: Towards Surprising Video Comprehension
% }

\title{
  \includegraphics[width=0.044\textwidth]{figures/logo.png} % 插入logo图片
  FunQA: Towards Surprising Video Comprehension % 标题文字
}

% The \author macro works with any number of authors. There are two commands
% used to separate the names and addresses of multiple authors: \And and \AND.
%
% Using \And between authors leaves it to LaTeX to determine where to break the
% lines. Using \AND forces a line break at that point. So, if LaTeX puts 3 of 4
% authors names on the first line, and the last on the second line, try using
% \AND instead of \And before the third author name.

\author{%
  Binzhu Xie\thanks{indicates equal contribution.\quad\textsuperscript{\Letter} Corresponding author. Contact: \texttt{ziwei.liu@ntu.edu.sg}}$^{~~,\dagger}$, Sicheng Zhang$^{\ast,\dagger}$, Zitang Zhou$^{\ast,\dagger}$,\\
  \textbf{Bo Li}$^{\ddagger}$, \textbf{Yuanhan Zhang}$^{\ddagger}$, \textbf{Jack Hessel}$^{\mathsection}$, \textbf{Jingkang Yang}$^{\ddagger}$, \textbf{Ziwei Liu}$^{\ddagger}$\textsuperscript{\Letter}\\
  $^{\dagger}$ Beijing University of Posts and Telecommunications, Beijing, China\\
  $^{\ddagger}$ S-Lab, Nanyang Technological University, Singapore\\
  $^\mathsection$ The Allen Institute for AI, WA, USA\\
  \texttt{\url{https://github.com/Jingkang50/FunQA}}
}


\maketitle
\vspace{-1cm}
\begin{figure}[htbp]
\centering
\includegraphics[width=\linewidth]{figures/fig_main.pdf}
\vspace{-0.6cm}
\caption{\textbf{Overview of FunQA.} FunQA comprises three subsets of surprising videos: \textit{1) HumorQA, 2) CreativeQA,} and \textit{3) MagicQA.} Each subset is associated with three common tasks: \textit{1) counter-intuitive timestamp localization}, \textit{2) detailed video description}, and \textit{3) reasoning around counter-intuitiveness} (see \textbf{H1-3}, \textbf{C1-3}, and \textbf{M1-3}). Furthermore, we offer higher-level tasks tailored for each video type, such as \textit{attributing a fitting and vivid title} for HumorQA and CreativeQA (see \textbf{H4}, \textbf{C4}), etc.
}
\label{figure_main}
\end{figure}


\begin{abstract}
Surprising videos, \emph{e.g.}, funny clips, creative performances, or visual illusions, attract significant attention. Enjoyment of these videos is not simply a response to visual stimuli; rather, it hinges on the human capacity to understand (and appreciate) commonsense \emph{violations} depicted in these videos.
We introduce \textbf{FunQA}, a challenging video question answering (QA) dataset specifically designed to evaluate and enhance the depth of video reasoning based on counter-intuitive and fun videos. Unlike most video QA benchmarks which focus on less surprising contexts, e.g., cooking or instructional videos, FunQA covers three previously unexplored types of surprising videos: \emph{1) HumorQA}, \emph{2) CreativeQA}, and \emph{3) MagicQA}.
For each subset, we establish rigorous QA tasks designed to assess the model's capability in counter-intuitive timestamp localization, detailed video description, and reasoning around counter-intuitiveness.
We also pose higher-level tasks, such as attributing a fitting and vivid title to the video, and scoring the video creativity.
In total, the FunQA benchmark consists of 312K free-text QA pairs derived from 4.3K video clips, spanning a total of 24 video hours.
Extensive experiments with existing VideoQA models reveal significant performance gaps for the FunQA videos across spatial-temporal reasoning, visual-centered reasoning, and free-text generation.
 \end{abstract}


% \section{Introduction}
% It is not a funny thing for AI to let AI model understand the joke of funny video. Consider the humorous video in Figure \ref{figure_main}. The boy was looking at his phone normally, but then was accidentally squeezed all over his face with ketchup. The contrast before and after is hilarious. Human can easily realize what is unusual about this video, while to make AI realize that there is something counter-intuitive in this video and explain humor is a challenging task. It needs a model to understand that squeeze ketchup on the face is uncommon and unexpected.\\
% Besides humorous videos, there are many other kinds of videos that need complex reasoning ability to accurately understand the counter-intuitive features of videos. For example, in Japan's Kasou Taishou [quote], people disguise themselves to achieve various unexpected effects, which are highly creative. The surprising effect of magic videos [quote] is also due to its counter-intuitive nature. AI needs to understand where these videos are counter-factual to explain why they are amazing. In other words, if the model can locate, describe and explain the counter-intuitive points on these videos, this will be strong support for the model's reasoning performance.\\
% In the task of picture reasoning, the model usually completes spatial relationship reasoning and infers possible behaviors [] according to experience. Because the complexity of humor comes from the difference between reality and expectation [] as long as the juxtaposition and transformation of incongruous things [], our task needs a model to better understand both temporal and spatial information. In this benchmark, the reasoning for video requires the model to have the ability of Spatial-Temporal reasoning, that is, these task cannot be completed through the pictures of a frame in the video.\\
% In some classic VideoQA datasets, such as TVQA \cite{lei2018tvqa} and MovieQA \cite{mun2017marioqa}, the input includes not only video, but also audio and even subtitle and bounding box. For reasoning on these datasets, the understanding of audio is indispensable. Different from this, we hope to focus on the visual understanding ability of the model. Therefore, our video will not eliminate sound, but the reasoning of our video can be completed only by visual content. More datasets, such as Social-IQ [], also want to perform audio-independence reasoning, but we have deeper requirements for visual understanding of video.\\
% In this paper, we have collected 10k humorous videos with an average duration of 20s, 3k creative super-variety videos and 1k magic videos, and our data set also includes 40k annotation. These counter-intuitive videos challenge AI's reasoning ability.\\
% Next, we propose three main reasoning tasks on the FunQA dataset:
% (1) time stamp localization, where model needs to find out the time stamps of the beginning and end of counter-intuitive segments in video.
% (2) description generation, where model objectively describes the counter-intuitive fragments in the video.
% (3) explanation generation, where the model need to generate detailed explanations on why the video is counter-intuitive, just like in the explanation provided by annotators in Fig \ref{figure_main}.
% \\At the same time, for three different categories of arts, we designed different subtasks. For humorous videos, we set up a subtask to generate funny video titles, inspired by those interesting titles on social media. For Kasou Taishou videos, there are also subtasks that generate caption. In addition, in Kasou Taishou, each performance will have an official rating, and we use this rating as the fifth subtask of creative video: rating the creativity of the video. For magic video, we hope that the model can understand the principle of magic, so the subtask is to explain the realization of magic effect.\\
% Unlike most Video QA datasets' multi-choice and open-ended questions [], our task requires the model to generate a free-text answer without templates or word bags.\\
% Our experiments show that after finetuning on this dataset , the performance of the model on the unknown counter-intuitive video has improved, which means the model can "understand" the jokes in the video better.\\
% (Limitations) However, some models failed on xxx tasks...\\
% We also show the performance of the model on other tasks after FunQA adjustment.\\
% Finally, we hope that our work can become a practical and challenging benchmark, which explores the boundaries of video reasoning with its rich content, and inspires future work. Our benchmark can be found at xxxxxx[link].\\

\section{Introduction}
The charm of surprising videos, be they funny, creative, or filled with visual illusions, % always --- I guess videos have not always existed :D 
% been reflected
%a rich source of entertainment widely enjoyed by humanity.
% These videos possess a universal allure,
offer enjoyment and commands attention from viewers.
This type of media elicits \emph{positive surprise} \cite{noordewier2013valence},\footnote{c.f., \emph{negative} surprise, e.g., a surprising medical bill.}
a captivating emotion that stems not merely from perceiving surface-level visual stimuli, but rather, the innate ability of humans to understand and find delight in unexpected and counter-intuitive moments \cite{martin1983humour}. However, despite significant advancements in today's computer vision models, the question remains: can video models ``understand" the humor/creativity in surprising videos? %, decipher the humor, uncover the creativity, and appreciate the magic in the same way that we do?

Consider the humorous video depicted in Figure~\ref{figure_main} (left) as an illustrative example. We witness a man engrossed in his phone, sharing a meal with friends. Suddenly, one of his companions squeezes a generous amount of ketchup, which, instead of adorning the fries, splatters onto the man's face. The shock in his eyes, combined with his ketchup-covered visage, elicits laughter.\footnote{The hostility/superiority theory of humor posits that
humor can arise from claiming superiority over someone or something \cite{gruner1978understanding,billig2005laughter}; but alternate (more optimistic) theories of humor exist, \cite{attardo2008primer} offers a survey.} While humans effortlessly recognize this as an unusual (and potentially entertaining) event, % it poses a considerable challenge for computer vision to grasp its underlying humor.
% While models At the very least,
the reasoning required to holistically understand the scene is complex:
a model needs to % comprehend the context of the video, where
recognize that individuals were \textit{gathered to enjoy a meal together}, and discern that the comedic element arises from the \textit{ketchup intended for the fries ending up on the man's face instead}, and that the innocent expression of the unfortunate victim indicates no significant harm was caused.

While there have been some efforts to enhance computer vision models' performance in Video Question Answering (VideoQA), these works have primarily focused on the common, less surprising videos found in existing VideoQA datasets. Examples of commonly employed VideoQA datasets include YouCook2 \cite{zhou2018towards} which contains video clips from 2K cooking videos, Howto100M \cite{miech2019howto100m} which consists of only instructional videos.
While there exist video datasets that explore the humor in TV shows \cite{mustard, hasan-etal-2019-ur} and include tasks such as predicting laughter tracks \cite{Patro_2021_WACV}, these tasks often heavily rely on audio and narrative cues, with visual clues might playing a lesser role.
% TVQA \cite{lei2018tvqa} and other datasets \cite that cover TV shows may contain surprising scenarios, but the questions either mostly focused on literal factors, e.g., who is talking to who. 
% \jack{is this true?}
% These benchmarks % fall short when it comes to assessing a。 deeper level of human understanding,
% may partly % particularly in comprehending videos that are intriguing or captivating.
% As a result,
% Models trained on these datasets may not access ``unusual" videos, resulting in limited video reasoning capability.
%For example, datasets like TVQA \cite{lei2018tvqa} include factual questions based on TV episodes, while NExT-QA \cite{xiao2021nextqanext} involves describing and explaining temporal actions. 

To address this gap and evaluate computer vision models' ability to identify and understand visual commonsense violations in videos, we introduce \textbf{FunQA}, a comprehensive and high-quality VideoQA dataset comprising 4.3K surprising videos and 312K manually annotated free-text QA pairs. Our dataset consists of three subsets: \emph{1) HumorQA}, \emph{2) CreativeQA}, and \emph{3) MagicQA}. Each subset covers different sources and video contents, but the commonality lies in their surprising nature, e.g., the unexpected contrasts in humorous videos, the intriguing disguises in creative videos, and the seemingly impossible performances in magic videos. Our experiments suggest that these surprising videos require different types of reasoning than common videos, as existing VideoQA methods perform poorly on the corpus. With FunQA, we hope to provide a benchmark that covers the popular, important, and sophisticated genre of counter-intuitive/surprising videos. %, but also, is offers a fun challenge for existing models. %exudes a fascinating and pleasing charm.

%\jingkang{introduce 3 video types and why we include them}


In FunQA, we formulate three rigorous tasks to measure models' understanding of surprise:
\textit{1) Counter-intuitive timestamp localization}: a model must % This task involves
identify the specific time period within a video when an unexpected event takes place.
\textit{2) Detailed video description}: a model must generate coherent and objective descriptions of the video content, evaluating models' fundamental video understanding capabilities.
\textit{3) Counter-intuitiveness reasoning}: a model must generate concrete explanations of why the video is surprising. % which requires deep reasoning about the counter-intuitive events in the video.
These tasks progressively assess the model's ability to perceive, articulate, and reason about the counter-intuitive elements present in surprising videos. Additionally, we propose auxiliary tasks that pose higher-level challenges within the benchmark including assigning an appropriate and vivid title to the video, etc. 
%These challenges include assigning an appropriate and vivid title to the video and evaluating its creativity, where the model assesses the level of ingenuity and novelty displayed in the video content.
% Contrary to traditional video QA datasets that depend on multiple-choice and open-ended questions, our tasks require the model to generate free-text responses. We've been cautious in this journey to avoid content that could be counter-intuitive yet distasteful. The annotations we collect are of high quality and objectivity by employing cross-validation. For each art, we set up a subset to balance the different tastes and agreements people have with counter-intuitiveness.
To summarize our contributions: 

\textbf{1) New VideoQA Dataset:} We build a large-scale dataset \textbf{FunQA}, which complements the existing VideoQA dataset with intriguing videos.

\textbf{2) Novel and Challenging Tasks:} We design a number of novel tasks that allow the model to explore previously untouched problems, such as timestamp localization, and reasoning around counter-intuitiveness. These tasks push video reasoning beyond superficial descriptions, demanding deeper understanding and discernment.

\textbf{3) Comprehensive Evaluation:} We have done an extensive and comprehensive evaluation of cutting-edge baselines, giving the field an insight and future research direction. 

% \begin{itemize}
%     \item \textbf{New VideoQA Dataset:} We build a large-scale dataset \textbf{FunQA}, which complements the existing VideoQA dataset with intriguing videos.
%     \item \textbf{Novel and Challenging Tasks:} We design a number of novel tasks that allow the model to explore previously untouched problems, such as timestamp localization, and reasoning around counter-intuitiveness. These tasks push video reasoning beyond superficial descriptions, demanding deeper understanding and discernment. 
%     \item \textbf{Comprehensive Evaluation:} We have done an extensive and comprehensive evaluation of cutting-edge baselines, giving the field an insight and future research direction. 
% \end{itemize}
%Although FunQA enhances the model's performance on counter-intuitive videos, the journey to fully comprehending surprising content is still far from over. Significant performance gaps exist between model and human responses, showing models' limitations in spatial-temporal reasoning, visual-centered reasoning, and free-text generation. With FunQA, we aspire to provide a benchmark that pushes the limits of video reasoning, opens new avenues of exploration, and fuels progress in AI's capability to understand surprising videos. We believe that FunQA is a step forward towards endowing AI with a more human-like comprehension of the world. Our FunQA dataset and related resources are available at \href{https://github.com/Jingkang50/FunQA}{https://github.com/Jingkang50/FunQA}.

% In total, the FunQA benchmark is a challenging and high-quality VideoQA dataset specifically designed to evaluate and enhance the depth of video reasoning based on counter-intuitive and fun videos containing 312K free-text QA pairs spread across 24 video hours.\\
% We find that, while fine-tuning on FunQA does improve the model's performance on counter-intuitive videos, the journey to fully comprehending surprising content is a long one. There are significant performance gaps, especially in areas of spatial-temporal reasoning, visual-centered reasoning, and free-text generation.With FunQA, we hope to provide a benchmark that tests the boundaries of video reasoning, opens up new avenues of exploration, and sparks progress in AI's capability of understanding surprising videos. We believe that FunQA is a stepping stone towards making AI not just smarter, but also more human-like in its understanding of the world. Our FunQA dataset and related resources are available at \href{https://github.com/Jingkang50/FunQA}{https://github.com/Jingkang50/FunQA}. 
%Counter-intuitive experiences, those which subvert the norm, evoke a spectrum of emotions from curiosity and surprise to amusement or even annoyance. We've been cautious in this journey to avoid content that could be counter-intuitive yet distasteful. We guarantee that each video, despite being counter-intuitive, exudes a fascinating and pleasing charm. This aligns perfectly with our fundamental objective: for the AI model to comprehend these surprising videos and experience the same level of enjoyment as humans do.\\
\section{Related Work}

% \textbf{Image Question Answering.} Visual Question Answering \cite{VQA} requires a model to answer visual related questions via understanding both visual content and question semantics. Existing benchmarks for visual question answering \cite{goyal2017making, zhu2016visual7w, krishna2016visual} use images \cite{VQA, zhu2016visual7w, gan2017vqs} and 2types of visual comprehension questions. They have made significant progress in assessing the visual language comprehension of systems from multiple perspectives of perception. Beyond factual reasoning, more research has fallen on unusual pictures, such as the recent GPT-4 \cite{openai2023gpt4} which can interpret the meaning of funny images (An example is "Lightning Cable"),  \cite{hessel2022androids} devised benchmarks for \textit{New Yorker-caliber} cartoon. comprehension, and \cite{Whoop} which devised benchmarks for commonsense-defying images comprehension.



% Please add the following required packages to your document preamble:
% \usepackage{booktabs}
% \usepackage{multirow}
% \usepackage[table,xcdraw]{xcolor}
% If you use beamer only pass "xcolor=table" option, i.e. \documentclass[xcolor=table]{beamer}
\begin{table}[!t]
\caption{\textbf{Comparison between FunQA and other existing benchmarks.}
Compare to other datasets, FunQA revolves around the captivating realm of interesting and counter-intuitive videos. The tasks within FunQA are specifically designed to challenge the vision capabilities of models, requiring strong skills in producing an in-depth description, interpretation, and spatial-temporal reasoning.
% FunQA's domain is interesting and counter-intuitive videos, FunQA selects videos that are humanly and repeatedly screened for high quality, FunQA's tasks focus on vision, and FunQA's answers require strong description, interpretation, and spatial-temporal reasoning skills already powerful in Free-text generation. FunQA is therefore designed to look at deeper description and explanation. 
Here we clarify the abbreviation in the table. \textbf{Anno.}: Annotation; \textbf{M}: Manual, \textbf{A}: Automatic; For Input, \textbf{V, A, S,} and \textbf{B} denote Video, Audio, Subtitle, and Bounding-box; \textbf{VC} means visual-centric, \textbf{Desc.} means Description, \textbf{Expl} for Explanation, \textbf{STR} for Spatial-temporal Reasoning. For QA Tasks, \textbf{MC} denotes Multiple Choice QA, \textbf{OE} means Open Ended QA, and \textbf{FT} means Free Text QA.
\label{T:main}}
\vspace{-0.05mm}
\resizebox{1.0\linewidth}{!}{
\begin{tabular}{@{}
>{\columncolor[HTML]{FFFFFF}}l 
>{\columncolor[HTML]{FFFFFF}}l 
>{\columncolor[HTML]{FFFFFF}}c 
>{\columncolor[HTML]{FFFFFF}}c 
>{\columncolor[HTML]{FFFFFF}}c 
>{\columncolor[HTML]{FFFFFF}}c 
>{\columncolor[HTML]{FFFFFF}}c 
>{\columncolor[HTML]{FFFFFF}}c 
>{\columncolor[HTML]{FFFFFF}}c 
>{\columncolor[HTML]{FFFFFF}}c 
>{\columncolor[HTML]{FFFFFF}}c 
>{\columncolor[HTML]{FFFFFF}}c @{}}



\toprule
\cellcolor[HTML]{FFFFFF}     
& \cellcolor[HTML]{FFFFFF}                           
& \cellcolor[HTML]{FFFFFF} 
& \multicolumn{3}{c}{\cellcolor[HTML]{FFFFFF}Video} & \multicolumn{5}{c}{\cellcolor[HTML]{FFFFFF}Question Answer}                         \\ 

\cmidrule(l){4-6} \cmidrule(l){7-12} 

\multirow{-2}{*}{\cellcolor[HTML]{FFFFFF}Datasets} 
& \multirow{-2}{*}{\cellcolor[HTML]{FFFFFF}Domain}                                        
& \multirow{-2}{*}{\cellcolor[HTML]{FFFFFF}Anno.} 
& Avg length (s) & \# Clips (K) & Input                                                           
& \# QA pairs (K) & VC & Desc. & Expl. & STR & QA Task      

\\ \cmidrule(r){1-12}

MarioQA \cite{mun2017marioqa} 
& Games  & M & 3-6 & 188 & V, A & 188 & Yes & No & Yes & Yes & OE \\
TGIF-QA \cite{jang-IJCV-2019} 
& Social Media & M & 3.1 & 71.7 & V & 165.2 & Yes & Yes & No & Yes & MC \& OE \\
MovieQA \cite{tapaswi2016movieqa}                                            
& Movies & A & 202.7 & 6.77 & V, A, S & 6.4 & No & No & Yes & Yes & MC \\
CLEVRER \cite{yi2020clevrer}                                           
& Synthetic Video & M & 5 & 20 & V & 305 & Yes & No & Yes & Yes & OE \\
TVQA \cite{lei2018tvqa}    
& TV shows & M & 76.2 & 21.8 &V, A & 152.5 & No & No & Yes & No & MC \\
TVQA+ \cite{lei2019tvqa}  
& TV shows & M &7.2 & 4.2 & V, A, B & 29.4 & No & No & Yes & Yes & OE \\
Social-IQ \cite{zadeh2019social}                                          
& Web videos & M & 99 & 1.25 & V, A, S, B & 7.5 & Yes & No & Yes & No & MC \\
NExT-QA \cite{xiao2021nextqanext}                                            
& Daily life & M & 44 & 5.4 & V, A & 52 & Yes & Yes & Yes & Yes & MC \& OE \\
KnowIT-VQA \cite{garcia2020knowit}                                        
& TV shows & M & 60 & 12 & V, A, S & 24 & Yes & No & Yes & Yes & MC \\
AGQA \cite{grundemclaughlin2021agqa}                                               
& Social Media & A & 30 & 9.6 & V, A & 192 & Yes & No & Yes & Yes & OE \\
AVQA \cite{yang2022avqa}                                               
& Social Media & M & 60 & 9.3 & V, A & 57.3 & Yes & No & Yes & Yes & MC \\
STAR \cite{wu2021star}                                               
& Daily life & A & - & 22 & V & 60 & Yes & No & No & Yes & MC  \\
Env-QA \cite{Gao_2021_ICCV}                                             
& Egocentric Video & M & 20 & 23.3 & V & 85.1 & Yes & No & No & Yes & MC \\
FIBER \cite{castro2022fiber}                                              
& Daily life & M & 10 & 28 & V, A, S & 2 & No & Yes & No & Yes & OE \\
\midrule

HumorQA                                           
& Daily life & M & 7 & 1.8 & V, A & 141.3 & Yes & Yes & Yes & Yes & OE \& FT \\

CreativeQA                                     
& Performance & M & 48 & 0.9 & V, A & 78.7 & Yes & Yes & Yes & Yes & OE \& FT \\

MagicQA                                  
& Magic & M & 10 & 1.6 & V, A & 91.9 & Yes & Yes & Yes & Yes & OE \& FT \\
\midrule
\textbf{FunQA} 
& \textbf{Surprising videos} 
& M & 19 & 4.3 & V, A & \textbf{312} & Yes & Yes & Yes & Yes & \textbf{OE \& FT} \\ \bottomrule 
\end{tabular}}
\vspace{-0.5cm}
\end{table}

\noindent\textbf{Video Question Answering Benchmarks}\quad
While the visual question answering (VQA) task focuses on enhancing models' ability in image comprehension \cite{goyal2017making, zhu2016visual7w, krishna2016visual}, video question answering (VideoQA) shifts the attention towards video comprehension. VideoQA is generally more challenging than VQA as it requires a comprehensive understanding of visual content, utilization of temporal and spatial information, and exploration of relationships between recognized objects and activities~\cite{yi2020clevrer}. To address the VideoQA task, the research community has introduced various benchmarks. As depicted in Table \ref{T:main}, Most commonly used VideoQA datasets are sourced from human-centric videos like movies \cite{tapaswi2016movieqa}, TV shows \cite{lei2018tvqa, lei2019tvqa, garcia2020knowit}, and social media \cite{jang-IJCV-2019, zadeh2019social, xiao2021nextqanext, grundemclaughlin2021agqa, yang2022avqa, wu2021star, castro2022fiber}, and there are also object-centric datasets of game videos \cite{mun2017marioqa}, synthetic videos \cite{yi2020clevrer} and egocentric videos~\cite{Gao_2021_ICCV}. MovieQA \cite{tapaswi2016movieqa} and TVQA \cite{lei2018tvqa} are commonly employed by VideoQA methods, which put forward tasks related to temporal and causal reasoning. However, they rely heavily on dialogue comprehension and textual plot summaries, which severely limits the challenge of visual reasoning. TGIF-QA \cite{jang-IJCV-2019} uses animated GIFs to challenge spatial-temporal reasoning, but as most GIFs are short videos of 3 seconds, and its tasks mainly focus on action description, TGIF-QA lacks complex reasoning evaluation ability. When most datasets use multiple choice questions as QA tasks, some methods, such as NExT-QA \cite{xiao2021nextqanext}, try to join open-ended questions. NExT-QA mainly focuses on daily life videos, but the open-ended answers are mostly simple sentences containing only a few words. To sum up, most existing methods focus on ordinary videos, lack of understanding of intriguing or unexpected videos, and advanced reasoning tasks such as generating complete explanatory texts of videos remain to be explored.\\
%As depicted in Table \ref{T:main}, there is a growing interest in utilizing question responses to measure video reasoning ability \cite{tapaswi2016movieqa, lei2018tvqa, lei2019tvqa, mun2017marioqa, jang-IJCV-2019, yi2020clevrer, xiao2021nextqanext, garcia2020knowit, grundemclaughlin2021agqa, wu2021star, Gao_2021_ICCV, castro2022fiber}. However, many well-known benchmarks generate data semi-automatically from popular TV shows or films \cite{tapaswi2016movieqa, lei2018tvqa, lei2019tvqa}, resulting in tasks and answers that rely heavily on dialogue and plot summaries rather than the visual content of the video \cite{tapaswi2016movieqa, lei2018tvqa,lei2019tvqa, kim2017deepstory, zadeh2019social}. This reliance on dialogue reduces the validity of the benchmarks for measuring visual reasoning \cite{tapaswi2016movieqa, lei2018tvqa, lei2019tvqa}.\\
% \textbf{Video Question Answering Benchmarks.} 
% As shown by the benchmarks in Table \ref{T:main}, there is growing interest in using question responses to measure video reasoning ability \cite{tapaswi2016movieqa, lei2018tvqa, lei2019tvqa, mun2017marioqa, jang-IJCV-2019, yi2020clevrer, xiao2021nextqanext, garcia2020knowit, grundemclaughlin2021agqa, wu2021star, Gao_2021_ICCV, castro2022fiber}. However, the data sources for these well-known benchmarks are often generated semi-automatically from well-known TV shows or films \cite{tapaswi2016movieqa, lei2018tvqa, lei2019tvqa}, so their tasks and answers often rely on dialogue and plot summaries rather than the visual content of the video \cite{tapaswi2016movieqa, lei2018tvqa, lei2019tvqa, kim2017deepstory, zadeh2019social}, leading to models that rely more on dialogue than on visual input, thus reducing the validity of the benchmarks for measuring visual reasoning \cite{tapaswi2016movieqa, lei2018tvqa, lei2019tvqa}.
% An increasing number of benchmarks are now often designed specifically to surprise by focusing on enhancing a particular domain of the model. \cite{yang2022avqa} focuses on improving the model's ability to process audio during problem solving. Moreover, \cite{xiao2021nextqanext, grundemclaughlin2021agqa, wu2021star}  currently focus on improving spatio-temporal reasoning.
\noindent\textbf{Video Question Answering Solutions}\quad
Earlier studies have explored various models, including LSTMs and graph-based neural networks, to capture cross-modal information \cite{li2019beyond, zhao2018open}. With the advent of Transformers, video understanding models like ClipBERT \cite{leiless} and CoMVT \cite{seo2021look} emerged, focusing on the comprehension of specific frames within a video. Subsequent models, such as Violet \cite{fu2021violet}, extended their ability to encompass temporal and spatial information. However, these methods have primarily been applied to short videos. In the realm of long videos, MIST \cite{gao2022mist} stands out by achieving state-of-the-art performance and excelling in terms of computation efficiency and interpretability. Furthermore, recent Vision Language Models (VLMs) such as \cite{li2023otter,2023videochat,maaz2023videochatgpt} have showcased remarkable video understanding capabilities.\\
\noindent\textbf{Counter-Intuitive Benchmarks}\quad
While many current computer vision benchmarks primarily focus on understanding commonsense content, there is a growing interest in addressing the realm of counter-intuitiveness. Several emerging benchmarks and models cater to this domain, such as Whoops~\cite{Whoop}, which emphasizes weird, unusual, and uncanny images, and MemeGraphs \cite{kougia2023memegraphs}, which revolves around memes featuring humor and sarcasm. Furthermore, some work even challenges models to comprehend complex multimodal humor in comics~\cite{hessel2022androids}. In the realm of large vision-language models, exemplified by GPT-4 \cite{openai2023gpt4}, there is a particular focus on showcasing their ability to provide explanations for funny pictures. However, when it comes to videos, existing datasets exploring humor in TV shows or comedy tend to heavily rely on audio and narrative cues \cite{mustard, hasan-etal-2019-ur,Patro_2021_WACV}, with visual clues playing a comparatively lesser role.
% However, there is no VideoQA benchmark involving counter-intuitive videos with deeper spatial-temporal reasoning.


% \textbf{Natural Text Generation} Natural language generation(NLG) is the task to automatically generate understandable texts, typically using a non-linguistic or textual representation of information as input \cite{covington2001building}. The first generation of metrics relied on handcrafted rules that measure the surface similarity between the sentences. To illustrate, BLEU \cite{papineni2002bleu} and ROUGE \cite{lin2004rouge}, two popular metrics, rely on N-gram overlap. Current approaches largely fall into two categories. Fully learned metrics, such as CIDEr \cite{vedantam2015cider}, BEER \cite{stanojevic2014beer}, RUSE \cite{shimanaka2018ruse}, and ESIM \cite{rebecq2018esim} are trained end-to-end, and they typically rely on handcrafted features and/or learned embeddings. Conversely, hybrid metrics,such as YiSi and BERTscore \cite{zhang2020bertscore} combine trained elements. Recently, there are four commonly used NLG benchmarks: all four systems employ the reference-based COMET \cite{rei2020comet},  BLEURT \cite{sellam2020bleurt}, and BERTSCORE \cite{zhang2020bertscore} metrics, in addition to the reference-free COMET-QE \cite{chimoto-bassett-2022-comet} metric.
% Recently, an increasing number of VQA tasks have now moved towards the use of open-ended answers (see in Table \ref{T:main}), which means that answers are no longer limited to predefined options. However, most answers are simply a few words or a fixed sentence. Few of the NLG that were generated in Large Language models were used as standard answers. \cite{hessel2022androids, Whoop} sets the standard answer to the QA task to a large section of natural language generated by LLM. In order to distinguish our approach from the prevalent practice of still confining answers within a certain range in open-ended VQA tasks, we have adopted a format where our annotated answers consist of lengthy paragraphs, aiming to provide as much coherent context as possible. We refer to this answer format as Free-text. 



\begin{comment}
Producing a good VQA dataset is complex, difficult and expensive because summarizing information from videos and giving inferred question-answer pairs is laborious, and most datasets will be produced with information from other modalities, such as subtitles and audio. \\
\textbf{General VQA datasets (factoid \& inference).}\quad Most of the early factoid VQA datasets were automatically generated by algorithms and were very homogeneous in form, which makes the learned
model easy to over-fit. They also in small dataset sizes and the length of videos are short, containing only a clip of the full video makes it difficult for the questions to involve complex reasoning.\cite{yu2019activitynet} For example, in VideoQA (FIB) \cite{zhu2017uncovering}, the main task is to describe the current clip and the questions are in the form of fill-in-the-blank questions, VideoQA \cite{zeng2017leveraging} uses a question generator to generate QA pairs from description of clips of movies, moreover, MSVD-QA, MSRVTT-Q \cite{xu2017video} and YouTube2Text-QA \cite{zhao2017video} generate questions of fixed categories in different ways, and their answers are all words.
Most inference VQA datasets ---xbz---. \\
So after that, some better datasets fall into several categories, including extremely large datasets, manually filtered datasets, and datasets for a particular domain.\\
\textbf{Large-scale VQA datasets.}\quad Among them, the availability of large datasets, for example, HowTo100M \cite{miech2019howto100m} and HowToVQA69M \cite{yang2021just} crawl Youtube videos in bulk and generate a large number of question-answer pairs using subtitles, Youtube-Temporal-180M \cite{zellers2021merlot}, a dataset for learning multimodal script knowledge, derived from 6 million public YouTube videos. All these databases provides the conditions for pre-training video understanding models.\\
\textbf{Manually annotated VQA datasets.}\quad In contrast, manually filtered and manually annotated datasets provide high quality datasets along with more difficult and complex tasks, providing the data base for training and evaluating models. For example, iVQA \cite{yang2021just}, which was manually selected from HowToVQA69M, eliminated QApairs that could be answered directly from the questions, and deliberately selected difficult questions, also ActivityNet-QA \cite{yu2019activitynet}, which is dedicated to generating better QApairs, fully annotated by humans, with more complex and inferential questions involving four categories: Motion, Spatial Relationship, Temporal Relationship and free.\\
\textbf{Specific VQA Datasets.}\quad Some datasets for specific types of videos provide help to solve specific tasks in a certain domain, such as the game dataset MarioQA \cite{mun2017marioqa}, the first-person dataset EgoVQA \cite{fan2019egovqa}, the wilderness dataset WildQA \cite{castro2022wild}, the chinese VQA dataset CREATE-210K and CREATE-3M/10M \cite{zhang2022create}.
\end{comment}

% \textbf{LLM based Models.} After the remarkable results of ChatGPT [20], GPT-4 [19] and other large language models [33, 32, 9], some studies have tried to include information from other modalities in the pre-trained language models.
\begin{comment}
\textbf{Transformer-based Works.}\quad Early transformer based video understanding models such as ClipBERT \cite{leiless} and CoMVT \cite{seo2021look} target the understanding of certain frames in the video, while later models such as Violet \cite{fu2021violet} add the understanding of temporal and spatial information.
\textbf{VLM.}\quad
Recently, large-scale models have been extensively studied. As a consequence, building on top of a powerful pretrained text-only language
model has been explored in concurrent work. VisualGPT (Chen et al., 2021) showed the benefit of
initializing the weights of a VLM with a pretrained language-only model for data-efficient training.
Tsimpoukelli et al. (2021) pushed this idea even further by freezing the pretrained LM weights. The
conditioning of the model on vision is similar to the idea of prefix tuning (Lester et al., 2021; Li and
Liang, 2021; Zhou et al., 2021; Zhu et al., 2021) (also called prompt tuning) where the prefix is
encoded by a trainable visual encoder. The benefit of building on top of a strong frozen language
model is that it may enable the visual language model to retain similar powerful language-only
abilities such as few-shot language adaptation, external knowledge retrieval, or dialogue capabilities.
Several works followed this idea with architectural differences in the conditioning of the frozen
language model. For instance, MAGMA (Eichenberg et al., 2021) adds bottleneck adapters (Houlsby
et al., 2019; Sung et al., 2021) within the frozen language model; ClipCap (Mokady et al., 2021)
proposes to use a vision-to-prefix transformer to map the vision features into a prefix instead of using
a simple linear layer mapping. VC-GPT (Luo et al., 2022) moves away from the visual prefix tuning
approach. They instead explore conditioning the frozen language model by grafting new learnt 
to the frozen language model. Finally, PICA (Yang et al., 2021) and Socratic Models (Zeng et al.,
2022) propose to use off-the-shelf vision-language models (Radford et al., 2021; Zhang et al., 2021)
to communicate the content of images using language descriptions to GPT-3 (Brown et al., 2020).
% new 3
Serveral recent models such as Flamingo \cite{alayrac2022flamingo}, LLaMA \cite{zhang2023llama} and AutoAD \cite{han2023autoad} are built based on large models such as Clip, GPT, etc., which demonstrate good results in the understanding of overall video information and adaptability to different complex problems. \\ 


\textbf{Fin-tune/pretrain on VQA models use special datasets.}\quad
Several recent methods pretrain multi-modal vision-language model. For example, --- \cite{} pretrain models using large amounts of narrated videos. In contrast to task-agnostic pretraining in the previous work, we show the benefits of task-specific pretraining for our target VideoQA task.
\end{comment}

% \textbf{Humor.} Attardo \cite{attardo2010linguistic} highlights three theories of humor: Hostility \cite{goldstein1976theoretical}, Release \cite{freud1960jokes,sep-humor}, and Incongruity.

% We focus on Incongruity that people laugh at the juxtaposition of incompatible concepts and at defiance of their expectations, that is, at the incongruity between expectations and reality \cite{kant1987critique,sep-humor}, and detecting dissonance between reality and expectations can place demands on the model's temporal capabilities. There has been some debate \cite{latta1999basic,boyd2004laughter,koestler2020act} attempting to clarify the roles of juxtaposition and shifting in humor. By designing humorous QA tasks, the model needs to detect juxtaposed and transformed events, corresponding to the two challenges of Multi-grained Visual Concepts and Multi-event proposed in MIST \cite{gao2022mist}, which also requires the use of spatial information

% \textbf{Creativity.} The standard definition is bipartite \cite{runco2012standard}: Creativity requires both originality and effectiveness. Mel Rhodes's answer to the question, "What is creativity?", is this: "The word creativity is a noun naming the phenomenon in which a person communicates a new concept." This paper is more focos on this originality and novelty than with effectiveness.

% \textbf{Magic.} The art of magic consists in creating illusions of the impossible \cite{nelms2012magic}, which contains a number of conjuring effects, such as appearance, vanish, and traspostition, these effects will be achieved using techniques, props or video clips. There are specific methods theory  
% \cite{lamont2005magic} to achieve each of conjuring effects, making seemingly counter-intuitive content interpretable.

% Incongruity in Humor, Originality and Novelty in Creativity and illusions of impossible in Magic are all \textbf{counter-intuitive} phenomena that require very strong priori-knowledge and spatio-temporal reasoning ability.
\section{FunQA Dataset}

% \begin{figure}[htbp]
%     \centering
%     \includegraphics[scale=0.5]{figures/dataset_main.pdf}
%     \caption{Pipleline of FunQA Dataset.}
%     \label{fig:dataset_main}
% \end{figure}

In this section, we provide a detailed explanation of the design principles that guided the creation of the FunQA dataset and its subsets. Additionally, we introduce our novel VideoQA tasks tailored for FunQA, and FunQA data statistics in Figure~\ref{fig:dataset statistics}. We introduce our construction pipeline in the end.


\begin{figure}[!t]

\begin{subfigure}[b]{0.3\textwidth}
    \centering
    \includegraphics[height=3cm]{figures/fig_word_cloud_h.png}
    \caption{Word cloud of HumorQA}
    \label{fig:wc_h}
  \end{subfigure}
  \hfill
  \begin{subfigure}[b]{0.3\textwidth}
    \centering
    \includegraphics[height=3cm]{figures/fig_word_cloud_c.png}
    \caption{Word cloud of CreativeQA}
    \label{fig:wc_c}
  \end{subfigure}
  \hfill
  \begin{subfigure}[b]{0.3\textwidth}
    \centering
    \includegraphics[height=3cm]{figures/fig_word_cloud_m.png}
    \caption{Word cloud of MagicQA}
    \label{fig:wc_m}
  \end{subfigure}

  \begin{subfigure}[b]{0.3\textwidth}
    \centering
    \includegraphics[height=3cm]{figures/fig_split_big.pdf}
    \caption{Number of videos and split}
    \label{fig:fun_split}
  \end{subfigure}
  \hfill  
 \begin{subfigure}[b]{0.3\textwidth}
    \centering
    \includegraphics[height=2.6cm]{figures/fig_video_len_big.pdf}
    \caption{Total video length}
    \label{fig:fun_len}
  \end{subfigure}
  \hfill
  \begin{subfigure}[b]{0.3\textwidth}
    \centering
    \includegraphics[height=3cm]{figures/fig_vqa_num_big.pdf}
    \caption{Number of QA pairs}
    \label{fig:fun_qa}
  \end{subfigure}
  \hfill


  
  \begin{subfigure}[b]{0.3\textwidth}
    \centering
    \includegraphics[height=3cm]{figures/fig_moment_distrib.png}
    \caption{Time Dist. of localization task}
    \label{fig:fun_moment}
  \end{subfigure}
  \hfill
  \begin{subfigure}[b]{0.35\textwidth}
    \centering
    \includegraphics[height=3cm]{figures/fig_avg_word_len.pdf}
    \caption{Avg. word count of answers}
    \label{fig:fun_avg}
  \end{subfigure}
  \hfill
  \begin{subfigure}[b]{0.3\textwidth}
    \centering
    \includegraphics[height=3cm]{figures/fig_val.pdf}
    \caption{Anno. consensus evaluation}
    \label{fig:fun_val}
  \end{subfigure}
  \hfill
  

  

  \caption{\textbf{Statistics of FunQA Dataset.} FunQA consists of three subsets, each corresponding to different video types, and is annotated with free-text QA pairs. The first row displays word clouds representing critical annotations for each subset. The second row provides key dataset statistics, including the number of videos for different splits, video length, and QA pair count for three subsets. In the last row, (g) highlights the high-frequency time span of the answer for localization questions in red, (h) shows the average word count of answers, and (i) presents the percentage of consensus between annotators for the same answer in a sampled set.
  }
  \vspace{-0.4cm}
  \label{fig:dataset statistics}
\end{figure}


\subsection{Video Selection}
% \noindent\textbf{Design principle}\quad 
In constructing the dataset, we adhered to three principles to address the challenges in video understanding capabilities: our dataset, FunQA, is \textbf{visual centered} and emphasizes \textbf{counter-intuitive reasoning}, \textbf{spatial-temporal reasoning}. Based on these principles, we collect 4365 videos from 3 different art genres and created three subsets: HumorQA, CreativeQA, and MagicQA.\\
% \jingkang{Based on the principle, we collect xxx videos from 3 different art genres....... What is your subset, why these datasets, simple stats of these dataset.}
\noindent\textbf{HumorQA}\quad
HumorQA composed of 1,769 meticulously curated web videos, serves as a unique source of insight into human humor comprehension. Notably, it contains the shortest average video length of 7s among the three subsets. We believe that the human process of understanding humor is complex and deep, requiring a holistic understanding of the video and adding a degree of common sense to it. Psychological research has demonstrated that humor arises from the incongruity \cite{kant1987critique,sep-humor} between reality and expectations, flourishing with the skillful juxtaposition and transformation of events \cite{latta1999basic,boyd2004laughter,koestler2020act}. This makes humorous videos a valuable asset for the VideoQA dataset, anticipated to enhance a model's proficiency in integrating information and performing deep reasoning.\\
\noindent\textbf{CreativeQA}\quad 
CreativeQA is a collection of 927 videos averaging 48s in length from a TV show called Kasou Taishou \cite{ntv}. This program, showcasing original and novel skits performed by various amateur groups and judged by a panel, boasts a strong creative flair 
 \cite{runco2012standard}. The essence of the show lies in using a mix of people and props to mimic reality, with audiences deriving pleasure from information integration and comparison. We anticipate that the imitation nature of the show will challenge the model's capacity for information extraction, while the longer video length and need for understanding creativity will put to test the model's comprehension of spatial-temporal information.\\
\noindent\textbf{MagicQA}\quad
MagicQA encapsulates 1672 magic performance videos sourced from across the web, spanning various genres like camera magic, close-up magic, and stage magic. The essence of magic revolves around the creation of seemingly impossible illusions \cite{nelms2012magic}, employing diverse effects such as disappearance, creation, and transformation. These illusions are infused with abundant spatial-temporal information. Through this dataset, we aim to empower the model to not only track the ensuing changes in objects but also unravel the underlying mechanics \cite{lamont2005magic} of these transformations.
% \noindent\textbf{FunQA}\quad
% In general, our meticulous filtering process ensures that the information within the aforementioned three types of art remains focused on visual content. These art forms naturally encompass complex spatial-temporal information, and their appreciation also requires deep human reasoning. Therefore, they align perfectly with our design guidelines.
\subsection{Task Definition}
% \jingkang{Each datasets has the following 4 types of tasks.....}


To comprehensively evaluate the model's ability to understand surprising videos, we designed the following 4 types of tasks for each subset: \\
\noindent\textbf{Counter-intuitive Timestamp Localization Task}\quad
The localization task is the base task to assess the model's comprehension abilities. It involves localizing counter-intuitive segments within the video, answers expressed in either seconds or frames. This task serves as the basis for the subsequent two main tasks in the three subsets, where the focus shifts to locating moments of humor, creativity, and magical effects, respectively. Successfully completing this task demands the model's understanding of the video's overall content, incorporating both temporal and spatial information.\\
\noindent\textbf{Detailed Description Task}\quad
The description task aims to evaluate the model's information extraction capabilities, serving as a fundamental aspect of video understanding. Across all three subsets, this task requires providing a free-text answer that describes the selected moment. Furthermore, this task allows for analysis of how the model extracts information and generates answers for subsequent tasks. By examining the model's performance in this task, we gain insights into its ability to extract relevant information and generate meaningful responses.\\
\noindent\textbf{Counter-intuitiveness Reasoning Task}\quad
The reasoning task is designed to test the model's ability to reason about the video, and in the three subsets, this question is Why Humorous, Why creative, and Why counter-intuitive and the answer is a free-text explanation. This task is very difficult and involves the model's deep reasoning ability; it requires the model to give a complete explanation using information from the entire video and its own common sense.\\
\noindent\textbf{Higher Level Tasks}\quad
In addition to the three main tasks, we design higher-level tasks to enhance the model's inference abilities on counter-intuitive videos. \textit{Title Task} in HumorQA and CreativeQA requires generating a concise title summarizing the video's content. \textit{Creative Scoring Task} in CreativeQA involves rating the creativity of videos between 0 and 20. \textit{Magic Method Task} in MagicQA requires the model to explain clearly the rationale behind the magic, and its purpose is to test the model's ability to reason more deeply. To ensure the accuracy of the answers, this task is only partially annotated and appears only in the test set, details of which can be found in Appendix \ref{appendix_a:dataset con}.
\\
In addition, some visual QA benchmarks \cite{Whoop} have begun to adopt free-text answers; this flexible answer format enables multimodal LLMs to go beyond multiple choice options (Large Language Models)~\cite{2023videochat,maaz2023videochatgpt}; for Video QA, however, free-text %while multiple-choice and open-ended formats continue to dominate the Video QA dataset.
are less common. To bridge this gap, we also consider some \textbf{free-text answer} evaluations, demanding enhanced video comprehension capabilities from the model. It will also assess the model's ability to generate reasonably lengthy textual responses.
% \input{table/construction}



\subsection{Dataset Statistics}
% FunQA contains \textbf{4,365} counter-intuitive video clips and \textbf{311,950} question-answer pairs, the total length of these videos is \textbf{23.9} h and the average length of video clips is \textbf{19} seconds. The distribution of the video lengths can be seen in Figure \ref{fig:fun_len}. These videos were divided into train / val / test set and the exact quantities can be seen in Figure \ref{fig:fun_split}. FunQA consists of three fine-grained sub-datasets, each sub-dataset contains four well-designed tasks, the distribution of specific tasks and their corresponding QApairs can be seen in Figure \ref{fig:fun_qa}. Figure \ref{} shows the word cloud map of FunQA as a whole, while Figures \ref{fig:wc_h}, \ref{fig:wc_c} and \ref{fig:wc_m} show the word cloud maps of the three sub-data sets, respectively., the other detailed statistics of sub-datasets are given in Table \ref{T:main} and Appendix \ref{section: FunQA statics}.\\ 
% FunQA has a well-established annotation process and high annotation quality, our cross-validation statistics can be seen in Figure \ref{fig:fun_val}, and the average length of our answers is up to -words, the specific data can be seen in Figure \ref{fig:fun_avg}.Besides, the heat map of the distribution of the results of the Counter-intuitive Timestamp Localization Task over the whole video can be seen in Figure \ref{fig:fun_moment}, which shows that the labeled results are in line with common sense.
% % % static table----------------------------
\begin{table}[!t]
\centering
\caption{\textbf{Statistics of the FunQA dataset.}
\label{T:statics}}
\resizebox{0.8\linewidth}{!}{

\begin{tabular}{lccccccc}
\hline
\multirow{2}{*}{Datase} & \multicolumn{4}{c}{Videos} & \multirow{2}{*}{Avg len(s)} & \multirow{2}{*}{Total len(h)} & \multirow{2}{*}{\# QA Pairs} \\ \cline{2-5}
           & Train & Val & Test & Total &    &       &       \\ \hline
HumorQA    & 1,244  & 350 & 172  & 1,769  & 7  & 3.72  & 141,250  \\
CreativeQA & 657   & 185 & 85   & 927   & 48 & 12.53 & 78,740  \\
MagicQA    & 1,171  & 334 & 167  & 1,672  & 16 & 7.65  & 91,960  \\ \hline
FunQA      & 3,072  & 869 & 424  & 4,365  & 19 & 23.90 & 311,950 \\ \hline \\
\end{tabular}
}
\end{table}
%static table ----------------------------------
% \\
% The distribution of question word length is shown in \ref{fig:fun_avg}(a). Questions are on average 34.24 words, which is much longer than existing VideoQA datasets (e.g., 8.7 in Activity-QA \cite{yu2019activitynet}, 11.6 in NExT-QA \cite{xiao2021nextqanext} ). And the median of each subtask is slightly less than the mean. The distribution of Moment is shown in \ref{fig:fun_moment}(b). We can clearly find that the distribution of counter-intuitive emotions in three different videos is very different. Humor, Creative and Magic videos account for 75\%, 23\% and 61\% of the videos respectively. HumorQA, Creative, Magic QA word cloud is shown in \ref{fig:wc_h, fig:wc_h, fig:wc_m}(c), FunQA full word clouds are in \ref{Data Statistics}.
FunQA contains \textbf{4,365} counter-intuitive video clips and \textbf{311,950} question-answer pairs, the total length of these videos is \textbf{23.9h} and the average length of video clips is \textbf{19} seconds. FunQA consists of three fine-grained subsets, each one containing well-designed tasks. The specific numbers and split of videos can be seen in Figure \ref{fig:dataset statistics} (d). Each subset's video lengths can be seen in Figure \ref{fig:dataset statistics} (d). The specific number of QA pairs for each task can be seen in Figure \ref{fig:dataset statistics} (e). \\
For our localization task, the timestamp heat map for the three different types of videos can be seen in Figure~\ref{fig:dataset statistics} (g), which shows the high-frequency time span of the answer. For the description and reasoning tasks, the average length of the words in their free-text answers reached \textbf{34.24}, which is much longer than existing VideoQA datasets (e.g., 8.7 in Activity-QA~\cite{yu2019activitynet}, 11.6 in NExT-QA~\cite{xiao2021nextqanext}). The specific word count of each task is shown in Figure \ref{fig:dataset statistics} (h). FunQA has a well-established annotation process and high annotation quality, the result of our annotation consensus evaluation can be seen in Figure \ref{fig:dataset statistics} (i). For each video category, more than 90\% of the annotations exhibit a high level of consensus, with only 1\% of the content showing low consensus. Approximately 8\% of the data demonstrates variations in consensus, thus highlighting the objectivity of our dataset. \\
HumorQA, CreativeQA, and MagicQA word clouds are shown in Figure \ref{fig:dataset statistics} (a-c). More statistics and FunQA full word cloud are given in Appendix \ref{appendix_a: Dataset statics}.

  
\subsection{Dataset Construction Pipeline} 
\label{Section: dataset construction pipeline}
FunQA dataset construction pipeline was in three stages: Pre-processing, Manual Annotation, and Post-Processing. The whole process took about 900 hours with over 50 highly educated undergraduates as part-time annotators, and we paid crowd-workers a target of ¥25/hr. More details of the dataset construction pipeline can be seen in Appendix \ref{appendix_a:dataset con}.\\
\noindent\textbf{Pre-processing}\quad
Initially, we crawled videos from TikTok, Bilibili, and YouTube (these videos are a collection of surprising videos). Then we performed a two-stage manual cleaning and cutting process on the collected videos to ensure counter-intuitive features and video quality and to exclude non-ethical content and sensitive information, resulting in a cut video clip.\\
\noindent\textbf{Manual Annotation}\quad
We annotated the videos according to the characteristics of different tasks design in Chinese. We screen and train the annotators to ensure the accuracy and high quality of the annotation, and finally produce the original annotated files. After the first round of annotation, we conducted a secondary annotation of 10\% of the tasks and performed Consensus Evaluation to ensure the objectivity of our annotations.\\
\noindent\textbf{Post-processing}\quad 
Based on our carefully designed tasks and high-quality annotations, we expanded our dataset using GPT-3.5. Firstly, we automatically translated the Chinese annotations into English. Subsequently, we generated more QA pairs that were faithful to the original ideas but presented differently. This not only made FunQA multilingual but also increased the number of QA pairs to 312K. Additionally, we created diverse task types, such as FunQA multiple-choice and FunQA dialogue. More details are given in Appendix \ref{appendix_a: Augmentation} and Appendix \ref{appendix_b: FunQA Extended dataset}.


% 

% Please add the following required packages to your document preamble:
% \usepackage{booktabs}
% \usepackage{multirow}
% \usepackage[table,xcdraw]{xcolor}
% If you use beamer only pass "xcolor=table" option, i.e. \documentclass[xcolor=table]{beamer}
\begin{table}[!t]
\caption{\textbf{Comparison between FunQA and other existing benchmarks.}
Compare to other datasets, FunQA revolves around the captivating realm of interesting and counter-intuitive videos. The tasks within FunQA are specifically designed to challenge the vision capabilities of models, requiring strong skills in producing an in-depth description, interpretation, and spatial-temporal reasoning.
% FunQA's domain is interesting and counter-intuitive videos, FunQA selects videos that are humanly and repeatedly screened for high quality, FunQA's tasks focus on vision, and FunQA's answers require strong description, interpretation, and spatial-temporal reasoning skills already powerful in Free-text generation. FunQA is therefore designed to look at deeper description and explanation. 
Here we clarify the abbreviation in the table. \textbf{Anno.}: Annotation; \textbf{M}: Manual, \textbf{A}: Automatic; For Input, \textbf{V, A, S,} and \textbf{B} denote Video, Audio, Subtitle, and Bounding-box; \textbf{VC} means visual-centric, \textbf{Desc.} means Description, \textbf{Expl} for Explanation, \textbf{STR} for Spatial-temporal Reasoning. For QA Tasks, \textbf{MC} denotes Multiple Choice QA, \textbf{OE} means Open Ended QA, and \textbf{FT} means Free Text QA.
\label{T:main}}
\vspace{-0.05mm}
\resizebox{1.0\linewidth}{!}{
\begin{tabular}{@{}
>{\columncolor[HTML]{FFFFFF}}l 
>{\columncolor[HTML]{FFFFFF}}l 
>{\columncolor[HTML]{FFFFFF}}c 
>{\columncolor[HTML]{FFFFFF}}c 
>{\columncolor[HTML]{FFFFFF}}c 
>{\columncolor[HTML]{FFFFFF}}c 
>{\columncolor[HTML]{FFFFFF}}c 
>{\columncolor[HTML]{FFFFFF}}c 
>{\columncolor[HTML]{FFFFFF}}c 
>{\columncolor[HTML]{FFFFFF}}c 
>{\columncolor[HTML]{FFFFFF}}c 
>{\columncolor[HTML]{FFFFFF}}c @{}}



\toprule
\cellcolor[HTML]{FFFFFF}     
& \cellcolor[HTML]{FFFFFF}                           
& \cellcolor[HTML]{FFFFFF} 
& \multicolumn{3}{c}{\cellcolor[HTML]{FFFFFF}Video} & \multicolumn{5}{c}{\cellcolor[HTML]{FFFFFF}Question Answer}                         \\ 

\cmidrule(l){4-6} \cmidrule(l){7-12} 

\multirow{-2}{*}{\cellcolor[HTML]{FFFFFF}Datasets} 
& \multirow{-2}{*}{\cellcolor[HTML]{FFFFFF}Domain}                                        
& \multirow{-2}{*}{\cellcolor[HTML]{FFFFFF}Anno.} 
& Avg length (s) & \# Clips (K) & Input                                                           
& \# QA pairs (K) & VC & Desc. & Expl. & STR & QA Task      

\\ \cmidrule(r){1-12}

MarioQA \cite{mun2017marioqa} 
& Games  & M & 3-6 & 188 & V, A & 188 & Yes & No & Yes & Yes & OE \\
TGIF-QA \cite{jang-IJCV-2019} 
& Social Media & M & 3.1 & 71.7 & V & 165.2 & Yes & Yes & No & Yes & MC \& OE \\
MovieQA \cite{tapaswi2016movieqa}                                            
& Movies & A & 202.7 & 6.77 & V, A, S & 6.4 & No & No & Yes & Yes & MC \\
CLEVRER \cite{yi2020clevrer}                                           
& Synthetic Video & M & 5 & 20 & V & 305 & Yes & No & Yes & Yes & OE \\
TVQA \cite{lei2018tvqa}    
& TV shows & M & 76.2 & 21.8 &V, A & 152.5 & No & No & Yes & No & MC \\
TVQA+ \cite{lei2019tvqa}  
& TV shows & M &7.2 & 4.2 & V, A, B & 29.4 & No & No & Yes & Yes & OE \\
Social-IQ \cite{zadeh2019social}                                          
& Web videos & M & 99 & 1.25 & V, A, S, B & 7.5 & Yes & No & Yes & No & MC \\
NExT-QA \cite{xiao2021nextqanext}                                            
& Daily life & M & 44 & 5.4 & V, A & 52 & Yes & Yes & Yes & Yes & MC \& OE \\
KnowIT-VQA \cite{garcia2020knowit}                                        
& TV shows & M & 60 & 12 & V, A, S & 24 & Yes & No & Yes & Yes & MC \\
AGQA \cite{grundemclaughlin2021agqa}                                               
& Social Media & A & 30 & 9.6 & V, A & 192 & Yes & No & Yes & Yes & OE \\
AVQA \cite{yang2022avqa}                                               
& Social Media & M & 60 & 9.3 & V, A & 57.3 & Yes & No & Yes & Yes & MC \\
STAR \cite{wu2021star}                                               
& Daily life & A & - & 22 & V & 60 & Yes & No & No & Yes & MC  \\
Env-QA \cite{Gao_2021_ICCV}                                             
& Egocentric Video & M & 20 & 23.3 & V & 85.1 & Yes & No & No & Yes & MC \\
FIBER \cite{castro2022fiber}                                              
& Daily life & M & 10 & 28 & V, A, S & 2 & No & Yes & No & Yes & OE \\
\midrule

HumorQA                                           
& Daily life & M & 7 & 1.8 & V, A & 141.3 & Yes & Yes & Yes & Yes & OE \& FT \\

CreativeQA                                     
& Performance & M & 48 & 0.9 & V, A & 78.7 & Yes & Yes & Yes & Yes & OE \& FT \\

MagicQA                                  
& Magic & M & 10 & 1.6 & V, A & 91.9 & Yes & Yes & Yes & Yes & OE \& FT \\
\midrule
\textbf{FunQA} 
& \textbf{Surprising videos} 
& M & 19 & 4.3 & V, A & \textbf{312} & Yes & Yes & Yes & Yes & \textbf{OE \& FT} \\ \bottomrule 
\end{tabular}}
\vspace{-0.5cm}
\end{table}


% \subsection{Dataset Comparision}

% As shown by the benchmarks in Table \ref{T:main}, there is growing interest in using question responses to measure video reasoning ability \cite{tapaswi2016movieqa, lei2018tvqa, lei2019tvqa, mun2017marioqa, jang-IJCV-2019, yi2020clevrer, xiao2021nextqanext, garcia2020knowit, grundemclaughlin2021agqa, wu2021star, Gao_2021_ICCV, castro2022fiber}. However, the data sources for these well-known benchmarks are often generated semi-automatically from well-known TV shows or films \cite{tapaswi2016movieqa, lei2018tvqa, lei2019tvqa}, so their tasks and answers often rely on dialogue and plot summaries rather than the visual content of the video \cite{tapaswi2016movieqa, lei2018tvqa, lei2019tvqa, kim2017deepstory, zadeh2019social}, leading to models that rely more on dialogue than on visual input, thus reducing the validity of the benchmarks for measuring visual reasoning \cite{tapaswi2016movieqa, lei2018tvqa, lei2019tvqa}.
% An increasing number of benchmarks are now often designed specifically to surprise by focusing on enhancing a particular domain of the model. \cite{yang2022avqa} focuses on improving the model's ability to process audio during problem solving. Moreover, \cite{xiao2021nextqanext, grundemclaughlin2021agqa, wu2021star}  currently focus on improving spatio-temporal reasoning.
% Other pure video benchmarks suffer from biases and simplicity associated with human-generated questions \cite{yu2019activitynetqa, tapaswi2016movieqa, jang-IJCV-2019, lei2018tvqa} or descriptions \cite{xu2017video, zeng2017leveraging}. The largest human annotated \cite{lei2018tvqa} and generated \cite{maharaj2017dataset} datasets contain 152.5K and 349K questions. Picture understanding has recently been increasingly studied in the context of strange, bizarre, and anomalous understanding, but there are no relevant datasets or benchmarks in the field of video understanding. FunQA surpasses the current VideoQA datasets in several aspects:\\
%  1) FunQA distinguishes itself through its diverse video types, each presenting unique characteristics within its respective subset. HumorQA involves incongruity and common sense elements, CreativeQA showcases intricate imitation behavior, and MagicQA entails the impossibility. FunQA's videos are longer and more informative, providing ample visual content for analysis.  \\
% 2) Secondly, FunQA places a strong emphasis on inference in its task design. This necessitates the model to possess deeper inference capabilities and leverage spatial and temporal information effectively. \\
% 3) Among the available VideoQA datasets, FunQA stands out for its inclusion of free-text type answers. Notably, the average answer length in Detailed Description Task and Counter-intuitiveness Reasoning Task is significantly longer compared to existing VideoQA datasets. This demands a higher output capability from the model and allows for evaluation using complex metrics. \\
% 4) FunQA benefits from high-quality human annotation, ensuring the reliability and accuracy of the dataset. They allow us to extend the FunQA dataset with GPT to suit different needs, e.g. FunQA Dialog is more in line with the LLM dialogue format, FunQA Multi-choice has proven evaluation metrics and testing is not limited to LLM, any model can be tested on it.



\section{Experiments}
% \jingkang{Please have 1-2 sentences to introduce what will you show in this section.}
In this section, we present an introduction to caption-based and instruction-based models, followed by an exploration of diverse metrics for evaluating FunQA tasks. Our comprehensive experiments and deep analysis of the results are then presented. More details are given in Appendix ~\ref{section_c: Experiments}.
\subsection{Baselines}
% \textbf{1) Q-ONLY} This baseline only uses the question feature as input to predict the answer.
% \textbf{2) V-ONLY} This baseline only uses the visual feature as input to predict the answer.
\subsubsection{Caption-based Models}
% \noindent\textbf{Caption-based Models:}\quad\\
\noindent\textbf{mPLUG}\quad
mPLUG \cite{li2022mplug} consists of two unimodal encoders for image and text independently, a cross-modal skip-connected network, and a decoder for text generation. Based on the connected representation of the image and prefix sub-sequence, the decoder is trained with a prefix language modeling loss by generating the remaining caption. \\
\noindent\textbf{GIT}\quad
GIT \cite{wang2022git} is composed of one image encoder and one text decoder, for videos, multiple frames are sampled and encoded independently, and features are added with an extra learnable temporal embedding before concatenation. The image encoder is based on the contrastive pre-trained model \cite{yuan2021florence} with raw image input and a compact 2D feature map output. The text decoder is a transformer module to predict the text description.
%which consists of multiple transformer blocks, each of which is composed of one self-attention layer and one feed-forward layer.


\subsubsection{Instruction-based Models}
\noindent\textbf{VideoChat}\quad
VideoChat \cite{2023videochat} is an end-to-end chat-centric video understanding system, the version we use is VideoChat-13B. Its VideoChat-Embed architecture is instantiated using BLIP-2 \cite{li2023blip2} and StableVicuna (13b-delta), and combines pre-trained ViT-G \cite{sun2023evaclip} and GMHRA \cite{wang2022internvideo}. For the token interface, VideoChat-Embed uses the pre-trained QFormer with an additional linear projection to output Video Embedding. Both video description and Video Embedding will be input in LLMs (LLAMA-13B \cite{touvron2023llama}) for multimodal understanding and output timestamped video text descriptions. 
%which has a VideoChat-Text architecture with video keyframes extracted using ffmpeg as input and a timestamped video text description as output
\\
\noindent\textbf{Video-ChatGPT}\quad 
Video-ChatGPT \cite{maaz2023videochatgpt} is a large vision-language model with a dedicated video-encoder and LLM, which feeds the video frames into pre-trained video encoder, adds spatio-temporal features and feeds them into linear layer. Video Embedding is then input into LLM (Vicuna-7B, v1.1) along with System command and User Query to output the answer. It uses a data-centric, human-assisted, and semi-automated annotation framework for high-quality video instructional data with unique multimodal (visual-verbal) capabilities.\\
\noindent\textbf{Otter}\quad 
The Otter model \cite{li2023otter} utilizes the OpenFlamingo training paradigm, where the pretrained OpenFlamingo model consists of a LLaMA-7B language encoder \cite{touvron2023llama} and a CLIP ViT-L/14 vision encoder \cite{radford2021learning}. In the fine-tuning process for instruction tuning, Otter freezes both the encoders and only fine-tunes the Perceiver resampler module. We evaluate two versions of Otter: one that is fine-tuned on the Dense Caption dataset \cite{krishna2017dense}, and another that is fine-tuned on the FunQA training set. Due to GPU memory limitations, Otter can only train and test on 128 frames from a video.

\begin{table}[!t]
\caption{\textbf{Main Results on FunQA Benchmark.} The FunQA benchmark consists of four task categories. H1, C1, M1 represent the counter-intuitive timestamp localization task, where \textbf{IOU} is used as the metric. H2, C2, M2 represent the detailed video description task, and H3, C3, M3 represent reasoning around counter-intuitiveness. For the higher-level tasks, H4, C4 involve attributing a fitting and vivid title. The responses for all these tasks in free-text format. We use the following metrics: \textbf{BLEU-4 / ROUGE-L / CIDEr} (shown in the first row) and \textbf{BLEURT / GPT-4} (shown in the second row) for evaluation. C5 represents scoring the video creativity, and the metric is the \textbf{Difference} between the predicted score and the official score. We tested the caption-based and instruction-based models. 
Here we clarify the abbreviation in the table.  \textbf{L.M.}: GIT\_LARGE\_MSRVTT; \textbf{L.V.}: GIT\_LARGE\_VATEX; \textbf{D.C.} means finetuned on Dense Caption; \textbf{FunQA} means finetuned on FunQA.\label{T:experiment}}
{\renewcommand\baselinestretch{1.5}\selectfont
% \Large
\setlength\tabcolsep{4pt} 
\resizebox{1.0\linewidth}{!}{
\begin{tabular}{lp{1cm}<{\centering} cccp{1cm}<{\centering}ccccp{1cm}<{\centering}cc}
\toprule
\multicolumn{1}{c}{\multirow{2}{*}{}} & \multicolumn{4}{c}{HumorQA} & \multicolumn{5}{c}{CreativeQA} & \multicolumn{3}{c}{MagicQA} \\

\cmidrule(l){2-5} \cmidrule(l){6-10} \cmidrule(l){11-13}
\multicolumn{1}{c}{} & H1 & H2 & \multicolumn{1}{c}{H3} & \multicolumn{1}{c}{H4} & C1 & C2 & \multicolumn{1}{c}{C3} & C4 & \multicolumn{1}{c}{C5} & M1 & M2 & M3 \\
\toprule
\rowcolor{TRAD_EVAL}
\multicolumn{13}{l}{\textbf{- Caption-based Model}} \\
mPLUG \cite{runco2012standard}  
& -
& \makecell[c]{\textbf{1.5} / 16.4 / \textbf{1.0} \\19.9 / 16.0 } 
& \makecell[c]{1.1 / 12.5 / 0.4 \\25.7 / 18.1 }  
& \makecell[c]{0.6 / 7.5 / 0.1 \\22.1 / 17.3 }  
& - 
& \makecell[c]{0.4 / 13.4 / 0.0 \\14.9 / \textbf{24.3} }  
& \makecell[c]{0.7 / 12.6 / 0.1 \\24.2 / 9.0 }  
& \makecell[c]{0.3 / 3.2 / 0.0 \\20.8 / 13.7 }
& \makecell[c]{-}  

& - 
& \makecell[c]{1.2 / 15.8 / 0.5 \\19.7 / 16.9 }
& \makecell[c]{0.9 / 8.9 / 0.4 \\21.2 / 8.8 }  \\
\midrule
% \rowcolor{TRAD_EVAL}
GIT (L.M.) \cite{wang2022git} 
& - 
& \makecell[c]{0.5 / 12.8 / 0.2 \\22.4 / 22.0 } 
& -
& \makecell[c]{\textbf{1.1} / 7.7 / \textbf{0.7} \\17.0 / 26.8 }

& - 
& \makecell[c]{0.0 / 6.4 / 0.0 \\14.4 / 5.0 } 
& \makecell[c]{-}  
& \makecell[c]{0.3 / 1.5 / 0.2 \\7.1 / 25.2 }
& -  

& -
& \makecell[c]{0.2 / 11.2 / 0.1 \\19.4 / 12.7 }
& -  \\
\midrule
% \rowcolor{TRAD_EVAL}
GIT (L.V.) \cite{wang2022git} 
& - 
& \makecell[c]{1.2 / 16.9 / 0.6 \\33.3 / 31.5  }
& - 
& \makecell[c]{1.0 / \textbf{8.8} / \textbf{0.7} \\\textbf{25.9} / 33.2}

& - 
& \makecell[c]{0.1 / 8.3 / 0.0 \\20.5 / 5.0 } 
& \makecell[c]{-}     
& \makecell[c]{\textbf{0.5} / 2.8 / \textbf{0.4} \\10.5 / 23.3 }
& \makecell[c]{-}

& - 
& \makecell[c]{0.6 / 13.7 / 0.1\\ 29.8 / \textbf{21.4} }  
& -  \\
\midrule

\rowcolor{NVL_EVAL}
\multicolumn{13}{l}{\textbf{- Instruction-based Model}} \\
VideoChat \cite{2023videochat} 
& - 
& \makecell[c]{0.5 / 13.7 / 0.0\\ \textbf{44.0} / \textbf{37.9} }
& \makecell[c]{0.5 / 13.5 / 0.0 \\ \textbf{45.4} / 31.9 }
& \makecell[c]{0.8 / 5.1 / 0.5 \\ 20.2 / \textbf{61.7} }
& - 
& \makecell[c]{0.3 / 7.5 / 0.0 \\ 21.7 / 10.9 }
& \makecell[c]{0.3 / 7.7 / 0.0 \\ 22.8 / 27.7 }
& \makecell[c]{0.2 / 1.2 / 0.2 \\ 7.3 / \textbf{51.1} }
& \makecell[c]{67.5}  
& - 
& \makecell[c]{0.6 / 15.5 / 0.0 \\ 47.4 / 14.2 } 
& \makecell[c]{0.3 / 9.2 / 0.0 \\ 43.1 / 24.6 }\\

\midrule
% \rowcolor{NVL_EVAL}
Video-ChatGPT \cite{maaz2023videochatgpt} 
& - 
& \makecell[c]{0.5 / 14.0 / 0.1 \\ 39.9 / 20.7 }
& \makecell[c]{0.7 / 12.4 / 0.1 \\ 40.1 / \textbf{33.0} }
& \makecell[c]{0.4 / 3.2 / 0.2 \\ 18.6 / 47.5 }
& - 
& \makecell[c]{1.1 / \textbf{19.8} / 0.2 \\ \textbf{45.8} / 19.1 }
& \makecell[c]{0.8 / 17.3 / 0.1 \\ \textbf{45.2} / \textbf{30.1} }
& \makecell[c]{0.2 / 1.9 / 0.2 \\ 18.8 / 44.5 }
& \makecell[c]{\textbf{85.4}}  
& - 
& \makecell[c]{0.7 / 20.8 / 0.0 \\ \textbf{50.0} / 11.8} 
& \makecell[c]{0.5 / 11.3 / 0.0 \\ 43.3 / \textbf{29.2} }\\
\midrule
% \rowcolor{NVL_EVAL}
Otter (D.C.) \cite{li2023otter} 
& - 
& \makecell[c]{1.1 / 14.3 / 0.4 \\ 30.2 / 9.8 } 
& \makecell[c]{1.2 / 14.2 / 0.4 \\ 32.3 / 13.9 } 
& \makecell[c]{0.5 / 5.4 / 0.1 \\ 21.7 / 13.3} 
& - 
& \makecell[c]{0.5 / 13.8 / 0.1 \\ 28.7 / 11.0 }  
& \makecell[c]{1.0 / 16.8 / 0.2 \\ 32.9 / 10.6} 
& \makecell[c]{0.3 / 2.3 / 0.1 \\ 17.7 / 4.2 }    
& \makecell[c]{45.0}  
& - 
& \makecell[c]{1.0 / 15.0 / 0.3 \\ 32.5 / 14.4 }  
& \makecell[c]{1.1 / 12.8 / 0.2 \\ 27.3 / 13.7}   \\
\midrule
% \rowcolor{NVL_EVAL}
% Otter (F.D.) \cite{li2023otter} 
% & - 
% & \makecell[c]{XX / XX / XX \\ XX / XX / XX}
% & \makecell[c]{XX / XX / XX \\ XX / XX / XX}
% & \makecell[c]{XX / XX / XX \\ XX / XX / XX}
% & - 
% & \makecell[c]{XX / XX / XX \\ XX / XX / XX}
% & \makecell[c]{XX / XX / XX \\ XX / XX / XX}
% & \makecell[c]{XX / XX / XX \\ XX / XX / XX}
% & \makecell[c]{XX}  
% & - 
% & \makecell[c]{XX / XX / XX \\ XX / XX / XX} 
% & \makecell[c]{XX / XX / XX \\ XX / XX / XX}\\
% \midrule
% \rowcolor{NVL_EVAL}
Otter (FunQA) \cite{li2023otter} 
& - 
& \makecell[c]{\textbf{1.5} / \textbf{18.1} / 0.9 \\ 38.4 / 22.2}  
& \makecell[c]{\textbf{1.3} / \textbf{15.4} / \textbf{0.5} \\ 42.6 / 31.0 }   
& \makecell[c]{0.8 / 5.9 /  0.5 \\ 24.5 /  51.3}   
& - 
& \makecell[c]{\textbf{1.5} / 19.6 / \textbf{0.5} \\ 40.0 / 21.9}   
& \makecell[c]{\textbf{2.2} / \textbf{21.2} / \textbf{0.5} \\ 41.1 / 21.1 }   
& \makecell[c]{0.3 / \textbf{4.3} / 0.3 \\ \textbf{21.7} / 43.9}    
& \makecell[c]{69.4}  
& - 
& \makecell[c]{\textbf{2.6} / \textbf{23.8} / \textbf{1.6} \\ 44.7 / 18.4 }  
& \makecell[c]{\textbf{3.4} / \textbf{20.3} / \textbf{2.6} \\ \textbf{44.5} / 19.8 }  \\
\bottomrule

\end{tabular}}\par
}
\vspace{-0.5cm}
\end{table}

\subsection{Evaluation Metrics}
% \textbf{Moment} 
% We use intersection of union to evaluate moment.\\
% \textbf{Score} The metric is $ (result - score)/20$\\
% \textbf{Free-text}
% 1) NLG. NLG traditional metrics such as BLUE \cite{papineni2002bleu}, ROUGE-1/2/L \cite{lin2004rouge} and METEOR \cite{banerjee2005meteor} are common for natural language understanding, which enables researchers to fairly compare their models to others. In GEM \cite{gehrmann2021gem}, people . However, most of the tasks in our dataset are free-text, which is different from the rest of the extant VQA datasets, and most of the existing models fail to yield the desired results, which are close to 0 when evaluated with the existing NLG metric after FunQA testing.\\
% 2) GPT 
% Some work \cite{} has demonstrated that GPT can be used as a metric for NLG with good results, so we introduced GPT to help with the free-text evaluation. We carefully designed the prompt to achieve the desired effect, and GPT will focus on the consistency of the details of the two texts and the understanding of humor, creativity and magic effects, as detailed in \ref{Metics}.
\noindent\textbf{Timestamp Localization (H1, C1, M1)}\quad 
We employ the intersection of union based on time span. \\
% which is calculated as $\min(result\_end, groud\_truth\_end) - \max(result\_start, groud\_truth\_start)/\min(result\_end, groud\_truth\_end) - \max(result\_start, groud\_truth\_start)$ 
\noindent\textbf{Description \& Reasoning (H2-4, C2-4, M2-3)}\quad
For all the \textbf{free-text} tasks, we employ three approaches for evaluation. Firstly, we utilize traditional NLG(Natural Language Generation) metrics. We use BLEU-4 \cite{papineni2002bleu}, ROUGE-L \cite{lin2004rouge}, CIDEr \cite{vedantam2015cider}, and BLEURT \cite{pu2021learning} as our metrics. The first two rely on N-gram overlap, which is only sensitive to lexical variations and cannot identify changes in sentence semantics or grammar. The latter two are reference-based evaluation metrics. Secondly, several works \cite{wang2023chatgpt, fu2023gptscore, kamalloo2023evaluating, chiang2023can} have shown promising results in utilizing GPT as a metric for NLG. Therefore, we introduce GPT-4 to assist in evaluating free-text similarity. We carefully design the prompts to make it possible to give objective ratings as much as possible like a human being. More details of GPT-4 prompts and evaluation criteria are provided in Appendix 
\ref{appendix_c: Evaluation Metrics}.\\
\noindent\textbf{Creative Scoring (C5)}\quad
The evaluation uses the formula: $Metrics = 100 \times \left(1 - \frac{{\lvert Predict - GT \rvert}}{{20}}\right)$.


\subsection{Results and Observations}
Our results are summarized in Table \ref{T:experiment}. As an illustration, the responses of different models on HumorQA can be seen in Figure \ref{fig:model comparsion}. Overall, the performance of the models on the FunQA dataset is generally unsatisfactory, and we have made several key findings:\\
\noindent\textbf{Timestamp localization task is the most challenging.}\quad
Caption-based models, due to their emphasis on captioning tasks, tend to provide descriptions of the entire video even when tasked with timestamp localization (refer to Appendix \ref{appendix_c: Failure of t1}). Conversely, instruction-based models, which are typically derived from image-based VLMs, focus on specific keyframes rather than considering the entire temporal space of the video.\\
\noindent\textbf{No clear winner across all tasks.}\quad
Caption-based models excel in providing detailed descriptions but struggle in tasks that require reasoning, resulting in a notable performance gap between description tasks (e.g., H2) and reasoning tasks (e.g., H3). On the other hand, instruction-based models demonstrate stronger reasoning capabilities but tend to underperform in description tasks. One possible explanation is that instruction-based models may generate excessive information in their answers, including a significant amount of incorrect information.\\
\noindent\textbf{Performance varies greatly across different video types.}\quad
Most models can get a relatively accurate answer in humor and magic videos but have difficulty answering the questions in creative videos. One possible reason is that humor and magic videos often depict daily life that models have encountered previously, whereas creative videos contain content that models have never seen before, rendering them unable to generate new ideas and resulting in irrelevant and erroneous answers.\\
\noindent\textbf{Insufficient evaluation metrics for free-text tasks.}\quad
Traditional metrics yield near-zero scores on free-text questions, as they solely focus on short textual similarity. While BLEURT scores are significantly higher, they still fall short in evaluating more complex similarities. Intuitively, GPT-4 is found to show preliminary capabilities in assessing free-text in deep understanding, which will be detailed in Appendix \ref{appendix_c: Evaluation Metrics}. However, there are still issues of instability, where the same content can receive different scores.\\
\noindent\textbf{Finetuned Otter performs well on traditional metrics but lags behind in GPT-4 score.}\quad
We finetuned Otter on Dense Caption and FunQA, and Otter (FunQA) shows obvious performance advantages over Otter (D.C.). While Otter performs better in traditional metrics like ROUGE-L compared with other instruction-based models, the GPT-4 score of Otter (FunQA) underperforms. One possible reason revealed is that the input of Otter is only 128 frames sampled from the video, which is insufficient for comprehensive reasoning. Besides, the discrepancy between Otter's scores on traditional metrics and GPT-4 matches our finding of the lack of evaluation metrics.

%it can be concluded that traditional metrics appear less suitable for more in-depth evaluation.

%The inadequate and unbalanced information and the lagging loos function may lead to ineffective training.

\begin{figure}[!t]
    \centering
    \includegraphics[width=\linewidth]{figures/model_comparsion.pdf}
    \vspace{-0.6cm}
    \caption{\textbf{Model responses on a HumorQA example.} Here shows the answers given by VideoChat~\cite{2023videochat}, Video-ChatGPT~\cite{maaz2023videochatgpt}, and Otter~\cite{li2023otter} on HumorQA video in Figure~\ref{figure_main}. On task \textbf{H2}, \textbf{H3}, VideoChat has the best performance. On task \textbf{H4}, Video-ChatGPT and Otter answer better, which is in line with our experiment result in Table \ref{T:experiment}. However, the answers from all models are still far from the ground truth. The descriptions of details and counter-intuitive explanations have numerous shortcomings. For example, Video-ChatGPT added incorrect details to the description, such as "wearing sunglasses", the humorous reason for "throwing ketchup" was wrongly interpreted by VideoChat as "knocking over the ketchup bottle", etc. }
    \label{fig:model comparsion}
    \vspace{-0.5cm}
\end{figure}
\section{Limitations and Future Work}
%This paper has two limitations. Firstly, FunQA's input is limited to raw videos, which restricts the model's ability to process and analyze the content effectively. Secondly, our experimental results reveal a relatively low overall score. Therefore, future work could involve performing additional annotation of the video, such as panoramic segmentation of images, to create a more extensive and comprehensive VideoQA dataset. Additionally, exploring the development of novel evaluation metrics is a promising research direction in the VideoQA field. Secondly,  exploring the development of novel evaluation metrics holds promise as a research avenue within the VideoQA field.\\
%In conclusion, as a new benchmark, FunQA paves the way for exploring interesting videos and delving into deeper human understanding. By tackling the aforementioned limitations, future research can further enhance the performance of models in video comprehension tasks.
This paper has two limitations. \textbf{1)}\quad Current FunQA dataset mainly includes video-level data and annotations, but denser annotations (e.g., like PVSG~\cite{yang2023panoptic}) can be developed to explore more possibilities of video reasoning. Examples include detailed spatial and temporal annotations, such as captions corresponding with specific time axes and annotations of object level. \textbf{2)}\quad The raw annotations are completed by the annotator in Chinese. In the process of translating into English, we first use GPT to polish and supplement Chinese annotations, making the text as thorough as possible, as short or incomplete text may result in misunderstanding. However, there may still be differences due to cultural diversities between the two languages.

In the future, we expect to expand our dataset with denser and more diverse annotations. Also, new metrics will be explored to better evaluate models' performance, especially in open-ended questions which lack in-depth metrics. Finally, we hope to provide directions for advancing models toward deeper video reasoning. Specific ideas include that Otter can sample more frames from the input video to improve causal reasoning capacity. Furthermore, in the future, when refining tuning datasets such as MIMIC-IT~\cite{li2023mimic} for visual language models, it may be beneficial to incorporate the FunQA training set. However, to ensure fair comparisons and prevent data leakage, it is advised that future research does not utilize the FunQA testing set.
% \section{Conclusion}
We introduce FunQA, a challenging VideoQA dataset of counter-intuitive and fun videos designed to explore the depth of video reasoning. We established a benchmark of four rigorous tasks and performed a comprehensive evaluation of frontier baselines, which were revealed to demonstrate a significant performance gap compared to manual annotation. Most baselines struggle especially in tasks requiring higher-level reasoning, such as counter-intuitive timestamp localization. FunQA dataset and novel benchmark tasks provide valuable insight and future research direction. We expect FunQA to underpin the next generation of VideoQA research that pushes video reasoning beyond superficial descriptions toward deeper understanding and discernment. 
\newpage
\begin{ack}
This study is supported by the Ministry of Education, Singapore, under its MOE AcRF Tier 2 (MOE-T2EP20221- 0012), NTU NAP, and under the RIE2020 Industry Alignment Fund – Industry Collaboration Projects (IAF-ICP) Funding Initiative, as well as cash and in-kind contribution from the industry partner(s).
\end{ack}
% \newpage
\bibliographystyle{unsrt}
\bibliography{ref}
\newpage
\appendix
% \noindent\textbf{Annotators.}\qquad
% All annotators are students from high level Chinese universities and we have tested their language skills to ensure that they can produce high quality annotations. During the annotation process, some of the annotators showed extremely outstanding linguistic skills, which also provided us with selection criteria for segmenting the dataset.\\
% \noindent\textbf{Pre-processing.}\qquad
% Before annotation, we manually cleaned the collected videos with two stages of clips. In the first stage, we cut all videos into content-independent clips, for example, we cut KT into individual performances. In the second stage, we first remove all videos that are unethical and contain sensitive information we cleaned the Humor videos to remove clips that do not meet ethical standards and have sensitive information, we do not operate on KT clips, only the score of each performance clip is recorded, and we subdivided the magic videos into clips to ensure that each clip includes only 1-3 coherent magic effects to avoid exceeding the model's understanding ability.\\
% \noindent\textbf{Tasks Annotation.}\qquad
% 1) We have developed detailed annotation documentation, and all annotators undergo rigorous training before performing the actual annotation. 
% 2) For the tasks in free-text, in Humor QA and MagicQA we suggested answers between 2-4 short sentences, in CreativeQA we would like a longer answer from the annotators. We asked the annotators to use concise and objective statements and avoid subjective adverbs and adjectives to ensure that the answers were in a similar language style. 
% 3) Annotators can question the counter-intuitive properties of video content or the generalizability of answers during the annotation process, and we will document such issues and address them in the next stage.
% 4) Cross-validation: We arranged another group of people who did not participate in the labeling to cross-validate the QA pair into three grades: completely wrong, ambiguous or too subjective, and no problem, and the statistical results can be seen in Figure \ref{}. because people's subjective understanding of humor, creativity and magic are bound to differ, a small number of disagreements is a normal result, we fixed the completely wrong QA pair and kept the rest.
% \noindent\textbf{Post-Processing.}\qquad
% After completing a round of annotation, we will check and filter the annotations. To ensure the quality of annotations, annotations with low quality or ambiguity will be returned to the task annotation stage for the next round of annotation. After all task annotations are completed with high quality, we will enter the data augmentation phase.


\section{More details of FunQA dataset} \label{appendix_a: FunQA more details}
\subsection{Dataset Construction Pipeline} \label{appendix_a:dataset con}

\begin{figure}[b]
    \centering
    \includegraphics[width=\linewidth]{figures_appendix/M4.pdf}
    \vspace{-0.6cm}
    \caption{\textbf{Example of Magic Method Task.} During the annotation process, we discovered that even as humans, it is difficult to fully understand the complete principles behind the implementation of magic tricks right from the beginning.}
    \vspace{-0.4cm}
    \label{fig:M4}
\end{figure}

\begin{figure}[!b]
    \centering
    \includegraphics[width=\linewidth]{figures_appendix/annotation_priciple.pdf}
    \vspace{-0.6cm}
    \caption{\textbf{FunQA Video Selection and Clip Principle.} We have a zero-tolerance policy regarding the inclusion of offensive content in our dataset. During the video sourcing process (video selection and video clipping stages), we ensure that such content is completely eliminated.}
    \vspace{-0.4cm}
    \label{fig:annotation-principle}
\end{figure}

\begin{figure}[!b]
    \centering
    \includegraphics[width=\linewidth]{figures_appendix/anno-error.pdf}
    \vspace{-0.6cm}
    \caption{\textbf{Examples of incorrect annotations and the suggested modifications.} Here we clarify the meanings of different cell colors. \textbf{Red}: Errors made by the annotator during the annotation process. \textbf{Yellow}: Types of annotation errors. \textbf{Green}: Reasons for the errors and suggested corrections.}
    \vspace{-0.4cm}
    \label{fig:annotation-error}
\end{figure}

\begin{figure}[!b]
    \centering
    \includegraphics[width=\linewidth]{figures_appendix/prompt-trans.pdf}
    \vspace{-0.6cm}
    \caption{\textbf{GPT prompt for Translation and augmentation.} Under our carefully designed prompt, GPT-3.5 can automatically generate English translations and five synonymous sentences from our high-quality Chinese annotations, thereby expanding FunQA dataset.}
    \vspace{-0.4cm}
    \label{fig:prompt-trans}
\end{figure}

\noindent\textbf{Pre-processing \& Qualification}\quad
For videos related to humor and magic, we downloaded them from different streaming platforms, mostly in the form of compilations. For Creative videos, we downloaded 26 episodes publicly available from Tokyo TV in Japan. We provided rigorous training to the annotators to ensure high-quality video clips in the final compilation. Annotators who successfully completed the Clip task according to the requirements are considered qualified and can proceed to the next stage of annotation.

\noindent\textbf{Training \& Annotation}\quad
We conducted systematic training for all annotators who passed the previous round of annotation, focusing on different tasks: \\
For the \textit{Counter-intuitive Timestamp Localization Task (H1, C1, and M1)}, the annotation format is a pair of numbers enclosed in square brackets, [xxxx, xxxx]. We asked the annotators to record the time intervals in which they felt pleasure (or amusement or shock) while watching the video. 

For the \textit{Task Detailed Description Task (H2, C2, and M2)}, we requested objective descriptions of what happened at [xxxx, xxxx], emphasizing a "what you see is what you get" approach. It is important to note that the annotations should only cover the selected time intervals and should not include subjective adverbs (such as vividly, vividly, or wildly). When describing characters or objects, be concise and add modifiers if there is ambiguity. An example of a poor label is "a man and a woman," while a relatively better example is "a man wearing a red hat and a woman wearing an apron." 

For the \textit{Counter-intuitiveness Reasoning Task (H3, C3, and M3)}, explain why the video is interesting in the context of the overall content. This part requires interpretive answers based on analysis, reasoning, and prior knowledge to explain why the video is counter-intuitive.

Specifically for the \textit{Magic Method Task (M4)}, we found that most annotators were not professional magicians, and even when watching instructional magic videos, it was challenging for them to provide complete and reasonable annotations for magic methods. Therefore, for M4, we only provided a small number of examples, which do not exist in the FunQA dataset. Figure \ref{fig:M4} shows an example of M4.

We conducted strict quality reviews in real-time for the annotations, including but not limited to checking if the annotations meet the requirements, if the selected videos contain sensitive content, and if the annotation texts contain sensitive content. The annotation principles (video selection and video clip, etc) can be seen in Figure~\ref{fig:annotation-principle}. Figure \ref{fig:annotation-error} also illustrates some errors encountered during the annotation process and the suggested modifications provided. After completing all the annotations, we conducted a consensus evaluation study to assess the objectivity of our annotations. We randomly selected 10\% of the videos and asked annotators to provide a consistency score (high consensus, medium consensus, low consensus) between their first and second annotations, considering all the previous annotations shown to them.

\noindent\textbf{Copyright \& License}\quad
We respect the copyright of each video. For content downloaded from streaming platforms, we obtained permission from the respective uploaders. For \textit{Kasou Taisho}, we submitted a fair use request to the official website of Tokyo TV in Japan.

We respect the personal identity information of everyone appearing in our videos and always strive to eliminate offensive content. By conducting strict reviews of the annotators and real-time quality checks during each annotation process, we ensure the absence of offensive content.
\subsection{Language Augmentation} \label{appendix_a: Augmentation}
For each QA pair, we first asked GPT-3.5 to translate the previous Chinese answers into English, and then let GPT-3.5 give five separate answers with the same meaning but different linguistic expressions. Afterward, we filtered out incorrect and incomplete generations. Figure \ref{fig:prompt-trans} shows the prompt we gave to GPT-3.5.

However, as mentioned in the main text in the limitation, the translation of GPT-3.5 will be ambiguous when the original text uses Chinese harmonics and some special Chinese meanings, and this problem mostly occurs in the Title Task, we screened such problems by hand and modified them, as shown in Figure \ref{fig:trans-err}.

\begin{figure}[htbp]
    \centering
    \includegraphics[width=\linewidth]{figures_appendix/tran_err.pdf}
    \vspace{-0.6cm}
    \caption{\textbf{Issues in translation caused by Chinese-English cultural differences.} In this example of translating and expanding a humorous video title, GPT-3.5 failed to understand the meaning of the original Chinese title, and we filtered out such data and made manual changes.}
    \vspace{-0.4cm}
    \label{fig:trans-err}
\end{figure}

\subsection{Data Statistics}\label{appendix_a: Dataset statics}
We collected raw videos from multiple video platforms, including short videos, long videos, and video clips, and the detailed data can be seen in Table \ref{T:raw_statics}. The word cloud shaped by FunQA logo of the all annotation word is shown in Figure \ref{fig:wc_funqa}.

% static table----------------------------
\begin{table}[h]
\caption{\textbf{Statistics of the FunQA raw data.} 
\label{T:raw_statics}}
\centering
\resizebox{1.0\linewidth}{!}{

\begin{tabular}{cccccc}
\hline
Raw Data          & Type                                & Source                   & Videos & Avg.len(s) & Total.len(h) \\ \hline
\multirow{3}{*}{Humor} & \multirow{2}{*}{Daily Life (Human)} & Bilibili                 & 351    & 182  & 15.8   \\
                       &                                     & TikTok                   & 1296   & 14      & 5.18    \\
                       & Nature (Animal)                     & Bilibili                 & 230    & 133  & 8.52    \\ \hline
Creative             & Performance                         & Kasou Taishou            & 26     & 6060  & 43.77   \\\hline
\multirow{2}{*}{MagicQA} & Close up Magic                      & \multirow{2}{*}{Youtube} & 765    & 96  & 20.40   \\ 
                       & Camera Magic                        &                          & 334    & 152  & 14.17   \\ \hline
FunQA                    & \textit{-}                          & -                        & 3002   & 129   & 107.87  \\ \hline\\
\end{tabular}

}
\end{table}
%static table ----------------------------------
\begin{figure}[htbp]
  \centering
  \includegraphics[width=0.5\textwidth]{figures/fig_word_cloud_funqa.png}
  \caption{\textbf{Word cloud of FunQA free-text answer.}}
  \label{fig:wc_funqa}
\end{figure}

\section{FunQA Extension Datasets}\label{appendix_b: FunQA Extended dataset}
Our main objective in designing these extended datasets is to leverage our high-quality annotated data and provide a rich and suitable data format for models. In addition, we also wanted to test the capability of GPT-3.5 and the quality of the dataset, and it turns out that GPT-3.5 can expand our data tens of times after providing high-quality free-text annotations, and our FunQA dataset is extremely scalable.
\begin{figure}[htbp]
    \centering
    \includegraphics[width=\linewidth]{figures_appendix/prompt-mc.pdf}
    \vspace{-0.6cm}
    \caption{\textbf{GPT prompt for FunQA Multi-choice.} We give GPT-3.5 prompt as system, and then enter the description and explanation of the video, and GPT-3.5 will output the data in the form of multiple-choice questions according to the prompt.}
    \vspace{-0.4cm}
    \label{fig:prompt-mcqa}
\end{figure}

\begin{figure}[htbp]
    \centering
    \includegraphics[width=\linewidth]{figures_appendix/prompt-dia.pdf}
    \vspace{-0.6cm}
    \caption{\textbf{GPT prompt for FunQA Dialog.} We give GPT prompt as system, and then input the description and explanation of the video, and GPT-3.5 will output the data in the form of dialogue according to the prompt.}
    \vspace{-0.4cm}
    \label{fig:prompt-dia}
\end{figure}
\subsection{FunQA Multi-choice Dataset} \label{appendix_b: FunQA MC}
FunQA Multi-choice Dataset is prepared to provide training and testing for arbitrary models, in this dataset our QA pairs are in the form of multiple choice, the answer is a word, phrase, or short sentence, and the type of questions are all descriptions. An example of the prompt given to GPT-3.5 and the data generated can be seen in Figure \ref{fig:prompt-mcqa}.
\subsection{FunQA Dialog Dataset}\label{appendix_b: FunQA Dia}
Most of the current LLMs are in the form of dialogues. To cater to their data input, we produced the FunQA Dialog dataset, in which we used GPT-3.5 to convert QA pairs into recursive dialogues with added context. In addition to this, we also let GPT-3.5 freely generate future dialogue development, which greatly expands the amount of information. An example of the prompt given to GPT-3.5 and the data generated can be seen in Figure \ref{fig:prompt-dia}.

\section{More Details of Experiment} \label{section_c: Experiments}
\subsection{Significance of New GPT-4 Based Metric} \label{appendix_c: Evaluation Metrics}
\noindent\textbf{GPT-4 prompt design}\quad 
For each of the three tasks, we designed three prompts for scoring. \textit{For Detailed Description Task (H2, C3, and M2)}, we designed the prompt in five areas: text length, text content variation, text detail variation, logical text description, and linguistic ability.
For \textit{Counter-intuitiveness Reasoning Task (H3, C3, and M3)}, we designed the prompt in six areas: expressiveness of language, the logic of response, the common sense of response, understanding of counter-intuition, differences in text detail, and length of text.
For \textit{Title Task (H4 and C4)}, we used the description, comprehension, and title of the manually annotated video as a reference to score the new title. Each Prompt can be seen in Figure \ref{fig:prompt-eval-title}, \ref{fig:prompt-eval-des}, and \ref{fig:prompt-eval-exp}.

\begin{figure}[htbp]
    \centering
    \includegraphics[width=\linewidth]{figures_appendix/prompt-eval-title.pdf}
    \vspace{-0.6cm}
    \caption{\textbf{GPT-4 prompt for Title Evaluation.} We give GPT prompt as system, and then input a description, explanation, and two titles, the first one is our annotation, and the second one is the output of the model, and GPT-4 will evaluate the similarity between these two texts according to the prompt's requirements.}
    \vspace{-0.4cm}
    \label{fig:prompt-eval-title}
\end{figure}
\begin{figure}[htbp]
    \centering
    \includegraphics[width=\linewidth]{figures_appendix/prompt-eval-des.pdf}
    \vspace{-0.6cm}
    \caption{\textbf{GPT-4 prompt for Description Evaluation.} We give GPT-4 prompt as system, and then input two descriptions, the first one is our annotation, and the second one is the output of the model, and GPT-4 will evaluate the similarity between these two texts according to the prompt's requirements.}
    \vspace{-0.4cm}
    \label{fig:prompt-eval-des}
\end{figure}
\begin{figure}[htbp]
    \centering
    \includegraphics[width=\linewidth]{figures_appendix/prompt-eval-exp.pdf}
    \vspace{-0.6cm}
    \caption{\textbf{GPT-4 prompt for Explanation Evaluation.} We give GPT prompt as system, and then input two explanations, the first one is our annotation, and the second one is the output of the model, and GPT-4 will evaluate the similarity between these two texts according to the prompt's requirements.}
    \vspace{-0.4cm}
    \label{fig:prompt-eval-exp}
\end{figure}

\noindent\textbf{Comparison between GPT-4 and Traditional Metrics}\quad
The principle of the traditional metrics is relatively simple. The traditional metric principles are as follows.\\
\noindent\textbf{BLEU-4}\quad BLEU, full name is Bilingual Evaluation Understudy, is a commonly used machine translation evaluation metric. It evaluates how good a machine translation result is by comparing how well it matches the N-gram of one or more human-translated reference results, which is a sequence of N consecutive words. BLEU-4, i.e., evaluates how well a combination of two words (i.e., a binary) matches. BLEU introduces a correction factor, Brevity Penalty (BP), to avoid this problem, which penalizes machine translation if the result is shorter than the reference translation. The BLEU score is the geometric mean of the individual N-gram accuracy multiplied by the shortness penalty. That is, the score of BLEU takes into account the precision and length of the translation result.\\
\noindent\textbf{ROUGE-L}\quad
ROUGE, known as Recall-Oriented Understudy for Gisting Evaluation, is a commonly used evaluation method for tasks such as automatic digesting and machine translation. ROUGE is mainly evaluated by comparing the overlap between the generated abstracts and the reference abstracts. Among them, ROUGE-L is an important variant of ROUGE, where L stands for Longest Common Subsequence (LCS), i.e., the longest common subsequence. Unlike the n-gram, the longest common subsequence does not require consecutive occurrences of items in the sequence.
\\
\noindent\textbf{CIDEr}\quad
CIDEr, known as Consensus-based Image Description Evaluation, is an evaluation metric specifically designed for evaluating image description (Image Captioning) tasks. The main advantage of CIDEr is that it can capture more detailed information because it uses TF-IDF weights to emphasize n-grams that occur frequently in manual annotation but are not common in all image descriptions.\\
\noindent\textbf{BLEURT}\quad
BLEURT, full name Bilingual Evaluation Understudy with Representations from Transformers, is an evaluation method based on the transformer model, specifically for evaluating the output quality of machine translation and natural language generation tasks. Unlike traditional evaluation metrics such as BLEU and ROUGE, BLEURT does not directly compare the n-gram match between generated text and reference text but uses pre-trained language models (e.g., BERT) to understand the semantic information of text. the advantage of BLEURT is that it can capture the deep semantic information of text, and can address some problems that are difficult to be handled by traditional evaluation metrics (e.g. synonym substitution, utterance rearrangement) giving reasonable evaluation.\\
It can be seen that the principles of the above traditional NLG metric are relatively simple and lack the ability to understand and evaluate the text with complex logic. In contrast, GPT-4, as an LLM, has a high ability to understand the text and can understand the text first and evaluate the similarity through prompt design.

\noindent\textbf{Instability of GPT-4 responses}\quad
During the experimental process using GPT-4 as the metric, we found that in a few cases, GPT-4 would provide different evaluations for the same content. In such cases, we would perform multiple evaluations and take the average. Figure \ref{fig:gpt-dif} shows an example of GPT-4 generating different responses to the same content.

\begin{figure}[htbp]
    \centering
    \includegraphics[width=\linewidth]{figures_appendix/GPT-4-dif.pdf}
    \vspace{-0.6cm}
    \caption{\textbf{GPT-4 generates different responses to the same content.} User Input is the two sentences input to GPT-4, which is used to judge the similarity. The GPT-4 Output 1 and GPT-4 Output 2 below are the answers given by GPT-4 in two different conversations, and you can see that the two assessments are different.}
    \vspace{-0.4cm}
    \label{fig:gpt-dif}
\end{figure}

\subsection{On Failure of the Timestamp Localization Task} \label{appendix_c: Failure of t1}
There is no model that has the ability to answer the Timestamp Localization Task. Figure \ref{fig:T1} shows exactly what is 
answered for each model.

\begin{figure}[htbp]
    \centering
    \includegraphics[width=\linewidth]{figures_appendix/Task_1.pdf}
    \vspace{-0.6cm}
    \caption{\textbf{Model responses on Timestamp Localization Task.} As can be seen, the model gives answers to the Timestamp Localization Task still focusing on describing the video content and still does not answer the specific time period when asked to type the number of frames and seconds.}
    \vspace{-0.4cm}
    \label{fig:T1}
\end{figure}

\subsection{Implementation Details} \label{appendix_c: Exp_result}
\noindent\textbf{mPLUG}\quad
mPLUG is a multi-modal system employing independent image and text encoders, a cross-modal network, and a text generation decoder, which is trained through prefix language modeling loss to generate captions from connected image and prefix sub-sequence representations.\\
\noindent\textbf{GIT}\quad 
GIT is a system with an image encoder and a text decoder; it processes multiple video frames independently, adds learnable temporal embeddings before concatenation, uses a contrastively pre-trained model for image encoding, and employs a transformer module for text prediction. We used the 14M version and used two models, GIT\_LARGE\_VATEX and GIT\_LARGE\_MSRVTT, which were fine-tuned on the video captioning task for the VATEX and MSRVTT datasets, respectively.\\
\noindent\textbf{VideoChat}\quad
VideoChat, specifically the VideoChat-13B version, is an end-to-end system for video comprehension that combines pre-trained models. It utilizes QFormer to generate video embeddings and then employs LLAMA-13B for multimodal understanding and outputs video text descriptions with timestamps. In the experiment, we used VideoChat-13B with the hyperparameters: beam search number = 1, temperature = 1, video segments = 8, and token = 512.\\
\noindent\textbf{Video-ChatGPT}\quad
Video-ChatGPT is a vision-language model with a video encoder and LLM. It generates answers using video embeddings and benefits from a data-centric, human-assisted annotation framework for high-quality video instructional data.
In the experiment, we used Video-ChatGPT-7B with its hyperparameter: temperature = 0.2, and token = 512.\\
\noindent\textbf{Otter}\quad
The Otter model employs the OpenFlamingo training paradigm, utilizing pre-trained encoders for language (LLaMA-7B) and vision (CLIP ViT-L/14). In the fine-tuning process, Otter prioritizes the Perceiver resampler module while keeping the encoders frozen. In the training stage, we finetuned Otter on Dense Caption and FunQA for a total of 3 epochs each. In the experiment, we used two versions of Otter with the same hyperparameters: beam search number = 3,  size of no-repeat-ngram = 0.2, and token = 256.

All of our experiments took five to six days on a single NVIDIA A40 48GB GPU.

\subsection{More Examples}
Figure \ref{fig:more-model-res} shows the responses of different models on CreativeQA and MagicQA.

\begin{figure}[htbp]
    \centering
    \includegraphics[width=\linewidth]{figures_appendix/More_Examples.pdf}
    \vspace{-0.6cm}
    \caption{\textbf{Model responses on CreativeQA and MagicQA. } For the description of the Creative video example, only VideoChat gives the key point of the bicycle, but its description also has many errors and omissions, and the remaining two models do not identify the bicycle. In the explanation task, the responses of all three models fail to clearly explain the creativity of this imitation performance. For the Magic video example, all three models perform very poorly in description and explanation, basically only answering the phone and the straw, but lacking the description and explanation of the magic effect.}
    \vspace{-0.4cm}
    \label{fig:more-model-res}
\end{figure}



\end{document}
